\chapter{理想流体动力学}

\begin{yijubox}
E1 \texttt{8-1考情分析.pptx} slide33：第七章理想流体动力学＝\textbf{概念为主}。\\
E2 \texttt{9-1考点分析.pptx}：\textbf{第六到第八章只有少量填空/判断/选择}，掌握基本概念＋公式套用即可。\\
E4 真题证据：2026 回忆版计算题第 4 题＝\textbf{证明势函数存在性}；
2015 年试题（科目 819）计算题第 7 题＝由流函数求速度分布与流线、第 8 题＝均匀流叠加点源求驻点与流线；
题库 \texttt{10.pdf} 选择第 13 题＝无旋运动的特征（涡量为零）。\\
\textbf{本章结论}：本章整体 \xstarfull{3}。第七章与第八章（黏性流体、边界层）是全书最"概念化"的两章：
大题只在极少数年份出现（2015、2026），但\textbf{一旦出现就落在「速度势函数与流函数」这一个考点上}；
其余内容（旋涡、空间势流）基本只考选择、判断、填空。\\
\textbf{复习策略}：① $\varphi$ 与 $\psi$ 的定义、存在条件、互求——必须能默写、能算；
② 几个基本势流的 $\varphi$、$\psi$、速度场——必须能默写；
③ 旋涡基本概念与四条定理——认得出、能判断对错即可。\\

\textbf{教材例题（第七章共 4 例，全部要精读）：}教材【例 7.1】（书 p121）\textbf{证明点源流动有势并求势函数}——就是 2026 回忆版计算题第 4 题的同型题；【例 7.2】（书 p123）由 $\psi=-\sqrt{3}\,x+y$ 求速度分布、流线方程与流量；【例 7.3】（书 p123）由例 7.1 的速度场求流函数（即 $\varphi$、$\psi$ 互求）；【例 7.4】（书 p146）证明均匀流的速度环量等于零。\\
\textbf{教材课后习题 7.1、7.2（书 p150）}的第 1、2 问就是“满足连续性方程否”“势函数、流函数存在否”，与 2026 年计算题第 4 题的问法完全一致。\\
\end{yijubox}

\section{速度势函数与流函数}

本章只讨论\textbf{不可压缩理想流体}的流动：无黏性（理想）$\to$ 不计切应力；不可压缩 $\to$ $\rho$ 为常数。

\subsection{平面流动与无旋流动}

\begin{kaodian}{5}{速度势函数 $\varphi$ 的定义与存在条件}
\textbf{平面流动：}流场中某一方向（如 $z$ 轴）流速为零（$u_z=0$），
且另两个方向流速 $u_x$、$u_y$ 与该方向坐标 $z$ 无关的流动，称为平面流动。

\textbf{无旋流动：}流体微团在运动中\textbf{不存在旋转运动}，即旋转角速度为零。
平面无旋流动 $w_z=0$，等价于 $\dfrac{\partial u_y}{\partial x}=\dfrac{\partial u_x}{\partial y}$。

\begin{gongshi}{速度势函数（势函数）$\varphi$ 的定义（教材式 7.1--7.3）}
\[ u_x=\frac{\partial\varphi}{\partial x},\qquad u_y=\frac{\partial\varphi}{\partial y},\qquad u_z=\frac{\partial\varphi}{\partial z}
   \qquad\Longleftrightarrow\qquad \bu=\grad\varphi \]
\textbf{含义：}无旋流动中三个速度分量可统一表示为同一个标量函数 $\varphi$ 的偏导数，
$\bu=\grad\varphi$。把 $\bu=\grad\varphi$ 代入无旋条件 $\bomega=\frac12\grad\times\bu=0$，
旋度自动为零——所以\textbf{有 $\varphi$ 就必无旋}，无旋就必有 $\varphi$，二者互为充要条件。\\
\textbf{适用条件：}\textbf{只要求无旋}（$\bomega=0$），\textbf{不要求}不可压缩、不要求恒定流动、
不要求是平面流；三维、可压缩流动同样可以有速度势。\\
\textbf{物理意义：}$-\grad\varphi$ 相当于一个"速度场"的势，故 $\varphi$ 称为速度势；
$\varphi=$ 常数 为\textbf{等势线（等势面）}，其法线方向就是速度方向。\\
\textbf{教材原话（书 p119）：}“由流动无旋条件确定的函数 $\varphi(x,y,z)$ 称为速度势函数，故无旋流又称为有势流，简称势流。”教材式 7.1 即 $\dd\varphi=u_x\dd x+u_y\dd y+u_z\dd z$，教材式 7.4 给出等势面方程 $\varphi(x,y,z)=C$。\\
\textbf{适用条件（教材口径，书 p120--121）：}教材把速度势函数的性质列为\textbf{四个特征}：①“不可压缩无旋流动的势函数是调和函数”（满足拉普拉斯方程）；②“存在势函数的流动是无旋流动”；③ 等势面与流线正交；④ 势函数的方向导数等于速度在该方向上的投影，即 $\partial\varphi/\partial s=u_s$。填空、简答按这四条写最稳。
\end{gongshi}

\begin{gongshi}{不可压缩时 $\varphi$ 满足拉普拉斯方程（教材 7.1.1 节 特征①，教材未给式号）}
\[ \grad\cdot\bu=\grad\cdot(\grad\varphi)=\grad^2\varphi=0 \qquad\Longrightarrow\qquad
   \frac{\partial^2\varphi}{\partial x^2}+\frac{\partial^2\varphi}{\partial y^2}+\frac{\partial^2\varphi}{\partial z^2}=0 \]
\textbf{含义：}不可压缩流动的连续性方程 $\grad\cdot\bu=0$，代入 $\bu=\grad\varphi$，
得到拉普拉斯方程 $\grad^2\varphi=0$；满足它的函数叫\textbf{调和函数}，
所以不可压缩无旋流动的 $\varphi$ 一定是调和函数。\\
\textbf{适用条件：}同时要求\textbf{无旋}（存在 $\varphi$）和\textbf{不可压缩}（$\grad\cdot\bu=0$）。
可压缩无旋流中 $\varphi$ 不满足拉普拉斯方程。\\
\textbf{教材原话（书 p120）：}教材把这一条列为速度势函数的\textbf{第 1 个特征}：“不可压缩无旋流动的势函数是调和函数”，并强调“\textbf{拉普拉斯方程的解具有可叠加性}”——这正是教材 7.3 节势流叠加原理的根据。
\end{gongshi}

\textbf{为什么要有势函数（教材引入背景）：}黏性流体绕物体流动时若雷诺数 $\Rey$ 较大，
黏性只体现在附面层内，附面层以外的流动可按无黏流动处理；
用势流理论求出附面层外边界上的压力分布与速度分布，作为附面层内求解的条件。
求机翼绕流升力时也把气流看作无黏流，用势流理论求解。因此势流理论是工程中绕流问题的理论基础。

\textbf{常考题型：}选择（"速度势函数存在于\rule{1.2cm}{0.4pt}流动中"，见下框真题）；\\
填空（"无旋流动的速度势满足\rule{1.8cm}{0.4pt}方程"$\to$ 拉普拉斯方程）；\\
计算（2026 回忆版计算题第 4 题，见后）——\textbf{证明某流动存在速度势}。
\textbf{重要度 \xstarfull{5}}\quad\yiju{E4 2026 回忆版计算题第 4 题"证明势函数存不存在"；E1 slide33"第七章＝概念为主"；E2 第六至第八章"掌握基本概念＋公式套用"}\quad\nandu{中}

\begin{zhenti}{2015 年试题 选择第 15 题}{单项选择题}
速度势函数存在于（\quad）流动中。\\
(A) 不可压缩流体\quad (B) 平面连续\quad (C) 处处无旋\quad (D) 任意平面
\end{zhenti}
\noindent\textbf{答案：(C)。}速度势存在的唯一条件是"处处无旋"，与是否可压缩、是否平面无关。

\begin{banfa}
\textbf{「证明某流动存在速度势」的三套等价打法（任选其一，考场上选最省力的）：}

\textbf{① 直接验证无旋（最常用）：}对平面流动只需算出 $w_z=\frac12\bigl(\frac{\partial u_y}{\partial x}-\frac{\partial u_x}{\partial y}\bigr)$，
证明它\textbf{恒为零}即可；三维流动则要验证 $w_x=w_y=w_z=0$。
算出的是"某一个含 $z$（或含其他坐标）的表达式"而不是 0，就说明不存在速度势。

\textbf{② 验证全微分：}证明 $u_x\dd x+u_y\dd y$（三维时再加 $u_z\dd z$）是某函数的全微分，
即验证交叉偏导相等 $\frac{\partial u_x}{\partial y}=\frac{\partial u_y}{\partial x}$——这与 ① 是同一个条件。

\textbf{③ 直接凑出 $\varphi$：}把 $u_x\dd x+u_y\dd y$ 沿路径积分
$\varphi=\int u_x\dd x+\int u_y\dd y$（第二步只补"含 $y$ 而不含 $x$ 的项"），
\textbf{只要凑得出 $\varphi$，存在性也就证完了}，而且顺手把答案拿到手。

\textbf{正反例速记：}
$u_x=ax$、$u_y=-ay$ 无旋（有 $\varphi$）；$u_x=-ay$、$u_y=ax$（绕原点刚性旋转）有旋
（此时 $w_z=a$，$\varphi$ 不存在）；$u_x=y$、$u_y=-x$ 亦有旋，无速度势。
\end{banfa}

\begin{zhenti}{教材【例 7.1】（书 p121）}{证明题}
一平面恒定不可压缩流动的流线为通过原点的向外发射的射线，速度大小 $u$ 反比于这点到原心的距离 $r$：$u=\dfrac{q}{2\pi r}$（$q$ 是正常数）。\textbf{证明这一流动是有势的}，求解势函数，并证明所得势函数是一调和函数。
\end{zhenti}
\noindent\textbf{教材解法（＝2026 回忆版计算题第 4 题的标准套路）：}
先写出速度分量 $u_x=\dfrac{q}{2\pi}\dfrac{x}{x^2+y^2}$、$u_y=\dfrac{q}{2\pi}\dfrac{y}{x^2+y^2}$；
再算 $\dfrac{\partial u_x}{\partial y}=\dfrac{\partial u_y}{\partial x}=-\dfrac{qxy}{2\pi(x^2+y^2)^2}$，两者相等即\textbf{有势}（$\varphi$ 存在）；
积分得 $\varphi=\dfrac{q}{2\pi}\ln\sqrt{x^2+y^2}$；最后算二阶偏导验证 $\dfrac{\partial^2\varphi}{\partial x^2}+\dfrac{\partial^2\varphi}{\partial y^2}=0$，故 $\varphi$ 是\textbf{调和函数}。
\\

\begin{tishi}
\textbf{教材同型例题与习题（专治“证明势函数存不存在”）：}\\
① 教材第七章只有 4 个例题，\textbf{【例 7.1】（书 p121）就是“证明流动有势并求势函数”}，与 2026 回忆版计算题第 4 题同型，务必精读；\\
② 教材\textbf{课后习题 7.1、7.2（书 p150）}的第 1、2 问是“满足连续性方程否”“势函数、流函数存在否”，问法与真题一致（习题 7.1：$u_x=4x+1$、$u_y=-4y$；习题 7.2：$u_x=x^2-y^2+x$、$u_y=-(2xy+y)$）；\\
③ 教材\textbf{没有}“给出速度场直接证明速度势不存在”的例题——凡题目问“不存在”，按上面的方法算出 $w_z\ne0$ 即可，判据仍是教材 7.1.1 节 的“存在势函数的流动是无旋流动”。
\end{tishi}
\end{kaodian}

\begin{kaodian}{5}{流函数 $\psi$ 的定义、存在条件与流量关系}
\textbf{流函数的引入：}由平面流动的连续性方程引入。
不可压缩平面流动的连续性方程为 $\dfrac{\partial u_x}{\partial x}+\dfrac{\partial u_y}{\partial y}=0$。

\begin{gongshi}{流函数 $\psi$ 的定义（教材式 7.5--7.8）}
\[ u_x=\frac{\partial\psi}{\partial y},\qquad\qquad u_y=-\frac{\partial\psi}{\partial x} \]
\textbf{含义：}形如 $\dd\psi=-u_y\dd x+u_x\dd y$ 的微分式在满足连续性方程时成为某函数 $\psi(x,y)$ 的全微分，
比较全微分定义即得上述两式。\textbf{流函数等于常数的曲线就是流线}。\\
\textbf{适用条件：}$\psi$ 只对\textbf{平面流动}存在（三维流动一般没有流函数）；
平面流动中\textbf{可压缩也可、无旋也可}。
对可压缩平面流动，定义推广为 $\rho u_x=\dfrac{\partial\psi}{\partial y}$、$\rho u_y=-\dfrac{\partial\psi}{\partial x}$
（此时 $\psi$ 的量纲不再是"面积／时间"，但 $\psi=$ 常数 仍是流线）。
轴对称流动\textbf{没有平面流函数}（上面这组定义式不适用），但\textbf{有轴对称流函数}——教材 7.5.2 节 专门给出柱坐标下的 $\psi(r,z)$（教材式 7.67、7.68），详见本章“空间势流”一节。
\end{gongshi}

\begin{tishi}
\textbf{（对旧稿的订正·轴对称流动有没有流函数）}本书旧稿此处写“轴对称流动没有流函数”，\textbf{与教材矛盾，现按教材订正}：\\
教材原话（书 p138）：“\textbf{轴对称流动是指流体在过某空间固定轴的所有平面上的运动情况完全相同的流动}”，并以“\textbf{轴对称流动的流函数}”为小节标题（教材 7.5.2 节）给出定义——\\
柱坐标（$r,\theta,z$）：$ru_r=\dfrac{\partial\psi}{\partial z}$、$-ru_z=\dfrac{\partial\psi}{\partial r}$（教材式 7.67、7.68）；\\
球坐标（$R,\theta,\beta$）：$\dfrac{\partial\psi}{\partial R}=-R\sin\theta\,u_\theta$、$\dfrac{\partial\psi}{\partial\theta}=R^2\sin\theta\,u_R$（教材式 7.69、7.70）。\\
正确说法是：\textbf{三维流动没有“平面流函数”，但轴对称流动有“轴对称流函数”}。轴对称流函数的两条性质与平面流函数完全平行：① 等流函数线就是流线；② 教材原话（书 p139）：“\textbf{在通过包含对称轴线的流动平面上，任意两点的流函数值之差的 $2\pi$ 倍，等于通过这两点间的任意连线的回转面的流量}”，即 $Q=2\pi(\psi_A-\psi_B)$。
\end{tishi}

\begin{gongshi}{$\psi=$ 常数 是流线}
\[ \dd\psi=-u_y\dd x+u_x\dd y=0 \qquad\Longrightarrow\qquad \frac{\dd x}{u_x}=\frac{\dd y}{u_y} \]
\textbf{含义：}$\psi$ 的等值线方程与流线微分方程完全相同，故 $\psi=c$ 就是一条流线。
这也是"流函数"这个名字的来源。\\
\textbf{适用条件：}平面流动；对不可压缩流动 $\psi$ 连续单值。
\end{gongshi}

\begin{gongshi}{流函数与流量的关系（教材 7.1.2 节 特性③，教材未给式号）}
\[ q_{\text{单宽}}=|\psi_2-\psi_1| \]
\textbf{含义：}\textbf{两条流线之间所通过的流量（单位宽度）等于这两条流线的流函数值之差的绝对值}。
二维平面问题取单位宽度，所以这是"单宽体积流量"，量纲 $\mathrm{L^2/T}$；
两条流线编号时按"$+y$ 侧减 $-y$ 侧"或者取绝对值，可避免符号纠结。\\
\textbf{适用条件：}不可压缩平面流动，且 $\psi$ 的定义式取 $u_x=\partial\psi/\partial y$、$u_y=-\partial\psi/\partial x$。
把每条流线的 $\psi$ 值减去参考流线的值，就得到参考流线以外的累计流量。\\
\textbf{教材原话（书 p121、p123）：}流函数的第 2 条特性是“\textbf{沿同一条流线的流函数是常数}”（令 $\psi=$ 一系列常数即得一系列流线）；第 3 条特性是“\textbf{两流线间单位厚度的流量等于其流函数值之差}”。填空就按这两句写。
\end{gongshi}

\begin{gongshi}{不可压缩平面流动的 $\psi$ 满足拉普拉斯方程}
\[ \grad^2\psi=\frac{\partial^2\psi}{\partial x^2}+\frac{\partial^2\psi}{\partial y^2}=0
   \qquad\text{（只在无旋时才成立）} \]
\textbf{含义：}若流动\textbf{同时}是平面、不可压缩、无旋（即平面势流），
由 $\frac{\partial u_x}{\partial y}=\frac{\partial u_y}{\partial x}$ 代入 $u_x=\psi_y$、$u_y=-\psi_x$ 即得 $\grad^2\psi=0$，
所以平面势流中 $\psi$ 也是调和函数。\\
\textbf{适用条件：}\textbf{无旋}是这里的硬条件；只不可压缩、不无旋（有旋）时 $\psi$ 依然存在，但\textbf{不}满足拉普拉斯方程。
\end{gongshi}

\begin{jielun}
\textbf{存在性一句话记牢（最易考、也最容易记反）：}\\
\textbf{流函数} $\psi$：只要\textbf{平面流动}就存在——可压缩、不可压缩、有旋、无旋全都行。\\
\textbf{速度势} $\varphi$：只有\textbf{无旋流动}才存在——平面可以、三维也可以，可压缩也可以，
但\textbf{只要有一点有旋，$\varphi$ 就不存在}。\\
所以：\textbf{有旋流有 $\psi$、没有 $\varphi$}；三维无旋流有 $\varphi$、没有 $\psi$（唯一例外是\textbf{轴对称流动}，它有教材 7.5.2 节 的轴对称流函数）。
\end{jielun}

\textbf{常考题型：}选择/判断（"有旋流动没有速度势但有流函数"；"流函数只存在平面流动中"）；
填空（"$\psi=$ 常数的曲线是\rule{1.6cm}{0.4pt}"$\to$ 流线）；
计算（2015 试题计算题第 7 题，见后）。
\textbf{重要度 \xstarfull{5}}\quad\yiju{E4 2015 年试题计算题第 7 题"由流函数求速度分布与流线"；E1 slide33"概念为主"；E2"掌握基本概念＋公式套用"}\quad\nandu{中}

\begin{zhenti}{2015 年试题 计算题第 7 题}{计算题（10 分）}
一平面定常流动的流函数为 $\psi(x,y)=-\sqrt{3}\,x+y$，
求速度分布，写出通过 $A(1,0)$ 和 $B(2,\sqrt{3}\,)$ 两点的流线方程。
\end{zhenti}
\noindent\textbf{解题要点：}$u_x=\dfrac{\partial\psi}{\partial y}=1$，$u_y=-\dfrac{\partial\psi}{\partial x}=\sqrt{3}$，
故速度场为均匀流 $\bu=(1,\sqrt{3}\,)$，大小为 2，方向与 $x$ 轴成 $60^\circ$。
流线就是 $\psi=$ 常数 的曲线：把 $A$、$B$ 的坐标代入得 $\psi_A=-\sqrt{3}$、$\psi_B=-\sqrt{3}$，
两点恰好在\textbf{同一条}流线上，故所求流线方程为 $y=\sqrt{3}\,x-\sqrt{3}$。

\noindent\textbf{（教材出处）这道题与教材\textbf{【例 7.2】（书 p123）完全相同}：}教材原题是“一平面恒定流动的流函数为 $\psi(x,y)=-\sqrt{3}\,x+y$，求速度分布，写出通过 $A(1,0)$ 和 $B(2,\sqrt{3})$ 两点的流线方程”，连数据都是同一组。
教材的解与结论：$u_x=\dfrac{\partial\psi}{\partial y}=1$、$u_y=-\dfrac{\partial\psi}{\partial x}=\sqrt{3}$，速度大小为 $2$、与 $x$ 正向夹角 $60^\circ$，教材原话“\textbf{这种速度分布不随地点变化的平面流叫平面均匀流}”；
$A$、$B$ 两点的流函数值均为 $-\sqrt{3}$，所以两点在同一流线上，流线方程为 $-\sqrt{3}\,x+y=-\sqrt{3}$（即 $y=\sqrt{3}\,x-\sqrt{3}$）；两点间流量 $\Delta\psi=0$（若追问“A、B 间流量是多少”，答“\textbf{同一流线上流量为零}”）。
另：教材\textbf{【例 7.3】（书 p123）}是“计算[例 7.1] 中平面流动的流函数”——$\varphi$、$\psi$ 互求的标准示范（把 $u_x$、$u_y$ 代入教材式 7.6 积分得 $\psi=\dfrac{q}{2\pi}\theta$）。

\begin{banfa}
\textbf{「由 $\psi$ 求速度与流线」三步：}
① 求速度：$u_x=\partial\psi/\partial y$、$u_y=-\partial\psi/\partial x$（\textbf{注意 $u_y$ 前有负号}）；
② 求流线：令 $\psi=c$，移项整理成 $y=(\cdots)x+(\cdots)$ 的标准式；
③ 定常数：把题给点坐标代入 $\psi$ 求 $c$——\textbf{若两个点算出同一个 $c$，说明两点在同一条流线上}，
题目只需写一个方程（2015 年这道题就是故意那样设计的，识别到这一点能省一半时间）。

\textbf{反向「由速度场求 $\psi$」：}用 $\psi=\int(-u_y)\dd x+\int u_x\dd y$，
第一步只需算出含 $x$ 的部分，第二步补上"不含 $x$ 的 $y$ 函数"，
最后用另一式（$u_x=\psi_y$ 或 $u_y=-\psi_x$）定出这个待定函数、校验答案。
\end{banfa}

\begin{yicuo}
\textbf{（本章第一号易错点）把 $\varphi$ 与 $\psi$ 的存在条件记反。}\\
$\varphi$ 只要\textbf{无旋}就存在（平面、三维、可压缩都行）；
$\psi$ 只要\textbf{平面}就存在（可压缩、有旋也行）。
"有旋流动没有速度势，但有流函数"——这句话是判断/选择题的标准答案。
反过来"三维流动有流函数"是\textbf{错}的，\textbf{轴对称流动却有轴对称流函数}（教材 7.5.2 节，书 p138），不要把这一条一起判错。
\end{yicuo}
\end{kaodian}

\begin{kaodian}{5}{速度势与流函数的关系：正交性、柯西-黎曼条件与互求}
\begin{gongshi}{等势线与流线正交（教材 7.1.3 节，教材未给式号）}
\[ \grad\varphi\cdot\grad\psi
   =\frac{\partial\varphi}{\partial x}\frac{\partial\psi}{\partial x}
   +\frac{\partial\varphi}{\partial y}\frac{\partial\psi}{\partial y}
   =u_x(-u_y)+u_y(u_x)=0 \]
\textbf{含义：}$\varphi$ 的梯度与 $\psi$ 的梯度互相垂直，所以\textbf{等 $\varphi$ 线与等 $\psi$ 线处处正交}——
等 $\varphi$ 线是等势线，等 $\psi$ 线是流线，两者构成正交网格，称为\textbf{流网}。\\
\textbf{适用条件：}平面不可压缩势流（同时存在 $\varphi$ 与 $\psi$）。\\
\textbf{教材原话（书 p124）：}“平面上若干等势线与流线构成了\textbf{正交曲线网}”——教材的术语是“正交曲线网”，本书其他地方叫“流网”，是同一个东西。
\end{gongshi}

\begin{gongshi}{柯西-黎曼条件（教材式 7.9）}
\[ \frac{\partial\varphi}{\partial x}=\frac{\partial\psi}{\partial y},\qquad\qquad
   \frac{\partial\varphi}{\partial y}=-\frac{\partial\psi}{\partial x} \]
\textbf{含义：}把 $u_x=\varphi_x=\psi_y$、$u_y=\varphi_y=-\psi_x$ 两式并写即得。
它是 $\varphi$ 与 $\psi$ 互相转换的\textbf{桥梁}：只要知道其中一个，就可以用这两个式子求出另一个。\\
\textbf{适用条件：}平面不可压缩势流。$-\psi_x$ 处的负号是这里的第二号易错点。
\end{gongshi}

\begin{banfa}
\textbf{互求两条路，考场上选简单的那条：}

\textbf{已知 $\varphi$ 求 $\psi$：}
① $u_x=\varphi_x$、$u_y=\varphi_y$；
② 用 $\dd\psi=-u_y\dd x+u_x\dd y$ 积分，
即 $\psi=\int(-u_y)\dd x+\int u_x\dd y$；
③ 校验：代回 $\psi_y$ 应等于 $u_x$（用 CR 条件收尾）。

\textbf{已知 $\psi$ 求 $\varphi$：}
① $u_x=\psi_y$、$u_y=-\psi_x$；
② 用 $\dd\varphi=u_x\dd x+u_y\dd y$ 积分；
③ 校验：代回 CR 条件。

\textbf{③ 三条自检口诀：}
"$\varphi_x$ 对 $\psi_y$、$\varphi_y$ 对 $-\psi_x$"；
"$\psi$ 的常数项 = $x$ 的待定函数，$\varphi$ 的常数项 = $y$ 的待定函数"；
结果里出现"$\ln r$"就应联想到源汇、出现"$\theta$"就应联想到环流、出现"$\cos\theta/r$"就应联想到偶极子——
\textbf{先看长相，能提前判断叠加了哪几个基本势流}。
\end{banfa}

\begin{yicuo}
\textbf{（本章第二号易错点）三个"线"分不清：}\\
$\psi=$ 常数 $\to$ \textbf{流线}；\quad $\varphi=$ 常数 $\to$ \textbf{等势线（不是流线）}；\quad
两者处处正交，正交网格叫\textbf{流网}。\\
另外：$\varphi$ 加上任意常数所描述的流场不变（因为 $\bu=\grad(\varphi+c)=\grad\varphi$），
$\psi$ 加上任意常数也不改变流场；
但"$\varphi$ 加常数"改变的是同一条等势线的取值，而"$\psi$ 加常数"改变的是流量起算基准——
\textbf{两者都不能改变速度场}，选择、判断题常拿这一点做文章。
\end{yicuo}

\textbf{常考题型：}计算（给出 $\varphi$ 求 $\psi$、或给出 $\psi$ 求 $\varphi$，并求速度场）；
填空（"等势线与流线\rule{1.4cm}{0.4pt}"$\to$ 正交；"柯西-黎曼条件为\rule{3cm}{0.4pt}"）；
判断（"$\varphi=$ 常数的曲线是流线"$\to$ 错，是等势线）。
\textbf{重要度 \xstarfull{5}}\quad\yiju{E4 2026 回忆版计算题第 4 题（势函数存在性）；E4 2015 年试题计算题第 7 题（由流函数求速度与流线）；E2"掌握基本概念＋公式套用"}\quad\nandu{中}
\end{kaodian}

\begin{kaodian}{2}{教材 7.1.4--7.1.5 节：复势、复速度与极坐标下的 $\varphi$、$\psi$}
\textbf{为什么补这一节：}教材 7.1 节专门讲了复变函数解法（选择题问过“平面流动的复势”“复速度的实部与虚部”），
而 7.2 节以后所有势流的 $\varphi$、$\psi$ 都是按极坐标给出的——这两处旧稿都没有，现按教材补上。

\begin{gongshi}{复势与复速度（教材式 7.10--7.12）}
\[ W(z)=\varphi+\mathrm{i}\psi,\qquad z=x+\mathrm{i}y,\qquad
   V=\frac{\dd W}{\dd z}=\frac{\partial\varphi}{\partial x}+\mathrm{i}\frac{\partial\psi}{\partial x}=u_x-\mathrm{i}u_y \]
\[ V=|u|\left(\cos\alpha-\mathrm{i}\sin\alpha\right)=|u|\mathrm{e}^{-\mathrm{i}\alpha} \]
\textbf{含义：}教材式 7.10 的 $W(z)$ 称为\textbf{复势}——把势函数与流函数拼成一个复变函数；
教材式 7.11 的 $V=\dd W/\dd z$ 称为\textbf{复速度}，其实部是 $u_x$、\textbf{虚部是 $-u_y$}；
教材式 7.12 说明 $V$ 的\textbf{模等于速度的大小} $|u|$、辐角等于速度方向角的相反数。
教材原话（书 p124）：“势函数和流函数是互为共轭的调和函数”“复速度的模等于速度的绝对值”
“可以根据共轭复变数的运算求出流场中每一点处的速度”。\\
\textbf{适用条件：}平面不可压缩势流（$\varphi$、$\psi$ 同时存在并满足柯西-黎曼条件）；
\textbf{虚数单位一律写 $\mathrm{i}$}（样式表没有专门的虚数单位宏）。\\
\textbf{常用复势对照：}均匀流 $W=u_\infty z$（教材式 7.18）；点源汇 $W=\pm\dfrac{q}{2\pi}\ln z$（教材式 7.28）；
点涡 $W=\dfrac{\Gamma}{2\pi\mathrm{i}}\ln z$（教材式 7.35）；圆柱绕流 $W=u_\infty z+\dfrac{M}{2\pi z}$（教材式 7.45）。
\end{gongshi}

\begin{gongshi}{极坐标下的 $\varphi$、$\psi$（教材式 7.13--7.15）}
\[ u_r=\frac{\partial\varphi}{\partial r}=\frac1r\frac{\partial\psi}{\partial\theta},\qquad\qquad
   u_\theta=\frac1r\frac{\partial\varphi}{\partial\theta}=-\frac{\partial\psi}{\partial r} \]
\[ \dd\varphi=u_r\dd r+r u_\theta\dd\theta,\qquad\qquad
   \dd\psi=-u_\theta\dd r+r u_r\dd\theta \]
\textbf{含义：}前两式是\textbf{极坐标下的柯西-黎曼条件}，把 $u_r$、$u_\theta$ 同时用 $\varphi$ 和 $\psi$ 表出；
教材式 7.14、7.15 是相应的积分式（对应直角坐标的 $\dd\varphi=u_x\dd x+u_y\dd y$、$\dd\psi=-u_y\dd x+u_x\dd y$）。\\
\textbf{适用条件：}平面不可压缩势流；$\theta$ 由 $+x$ 轴\textbf{逆时针}量起。\\
\textbf{用法：}题目给的是极坐标速度分量时（如 $u_r=U\cos\theta$、$u_\theta=-U\sin\theta$），
用 $\dd\psi=-u_\theta\dd r+r u_r\dd\theta$ 积一次就能得到 $\psi$，比换成直角坐标再算快得多。
\end{gongshi}

\textbf{常考题型：}选择/填空（“复势 $W(z)=\varphi+\mathrm{i}\psi$ 的实部是\rule{1.8cm}{0.4pt}、虚部是\rule{1.8cm}{0.4pt}”；
“复速度的模等于\rule{1.6cm}{0.4pt}”）；也可与 $\varphi$、$\psi$ 互求合并考小计算。
\textbf{重要度 \xstarfull{2}}\quad\yiju{E2“第六到第八章只有少量填空/判断/选择题，掌握基本概念＋公式套用即可”；E6 教材 7.1.4、7.1.5 节（书 p124--125）教材为独立小节}\quad\nandu{中}

\begin{banfa}
\textbf{① 复势三件套速记：}$W=\varphi+\mathrm{i}\psi$（实部势函数、虚部流函数）；
$V=\dfrac{\dd W}{\dd z}$（实部 $u_x$、\textbf{虚部 $-u_y$}，虚部带负号是唯一的坑）。\\
\textbf{② 极坐标看“管什么”：}$u_r$ 管进出、$u_\theta$ 管绕圈；求驻点就令两者同时为零。
\end{banfa}
\end{kaodian}

\section{几种基本平面势流}

以下四种流动是全部平面势流的"积木"。\textbf{每个都必须背出：$\varphi$、$\psi$、速度场、流线形状}。

\begin{kaodian}{4}{均匀直线流（等速均匀流）}
\begin{gongshi}{均匀直线流的 $\varphi$、$\psi$ 与速度场（教材式 7.16--7.19）}
\[ u_x=U,\quad u_y=0:\qquad \varphi=Ux,\qquad \psi=Uy \]
\[ \text{一般情形 } u_x=a,\ u_y=b:\qquad \varphi=ax+by,\qquad \psi=-bx+ay \]
\textbf{含义：}流场中各点速度矢量互相平行且大小相等。沿 $x$ 轴的均匀流（$U$ 沿 $+x$）
$\varphi=Ux$、$\psi=Uy$；沿 $y$ 轴的均匀流（$u_y=V$）则 $\varphi=Vy$、$\psi=-Vx$。\\
\textbf{适用条件：}无旋（显然 $w_z=\frac12(0-0)=0$）、不可压缩平面流动，
是全章所有叠加问题的"底料"。\\
\textbf{流线形状：}一组\textbf{平行于 $x$ 轴的直线}（$\psi=Uy=c$ 给出 $y=$ 常数）。\\
\textbf{教材原话（书 p125）：}“\textbf{平面均匀流指在同一时刻，流场中所有点的速度矢量的大小与方向都相同的平面流动}”；教材式 7.16--7.19 依次给出 $\varphi=u_\infty x$、$\psi=u_\infty y$、$W=u_\infty z$、$W=u_\infty z\mathrm{e}^{-\mathrm{i}\alpha}$，并指出\textbf{等势线是与 $y$ 轴平行的直线簇、流线是与 $x$ 轴平行的直线簇}。
\end{gongshi}
\textbf{常考题型：}叠加题的基本组成部分（2015 年计算题第 8 题＝"速度为 $v_\infty$ 的均匀流＋点源"）。
\textbf{重要度 \xstarfull{4}}\quad\yiju{E4 2015 年试题计算题第 8 题（均匀流叠加点源）；E2"掌握基本概念＋公式套用"}\quad\nandu{易}
\end{kaodian}

\begin{kaodian}{4}{点源与点汇}
\begin{gongshi}{点源的 $\varphi$、$\psi$ 与速度场（教材式 7.20--7.25）}
\[ u_r=\frac{Q}{2\pi r},\qquad u_\theta=0,\qquad \varphi=\frac{Q}{2\pi}\ln r,\qquad \psi=\frac{Q}{2\pi}\theta \]
直角坐标：$u_x=\dfrac{Q}{2\pi}\dfrac{x}{x^2+y^2}$，$u_y=\dfrac{Q}{2\pi}\dfrac{y}{x^2+y^2}$。\\
\textbf{含义：}流体从平面上一点 $o$ 沿径向均匀流向四周的流动称为\textbf{源流}，
$o$ 为\textbf{源点}；由源点流出的\textbf{单位厚度}流量 $Q$ 称为\textbf{源流强度}。
源点处 $r\to0$ 使 $u_r\to\infty$，故源点是\textbf{奇点}（流场中除源点外处处无旋）。
流速与半径 $r$ 成反比，恰使各同心圆周上的单宽流量都为 $Q$。\\
\textbf{适用条件：}不可压缩平面势流；$r>0$，源点本身无定义。\\
\textbf{流线形状：}以源点 $o$ 为起点的\textbf{半辐射线}（线族 $\theta=$ 常数）；
等势线是以源点为圆心的\textbf{同心圆}，与流线正交——正好构成流网。\\
\textbf{教材原话（书 p126）：}“\textbf{如果流体从某点向四周均匀径向流出}……称为点源”，“如果流体从四周往某点呈直线均匀径向流入……称为点汇”，“\textbf{流量 $q$ 为点源与点汇的强度}”。
注意：教材把源强记作 $q$、本书记作 $Q$，含义相同（单位厚度的体积流量）；教材式 7.20--7.25 依次给出 $u_r=\pm\dfrac{q}{2\pi r}$、$\varphi=\pm\dfrac{q}{2\pi}\ln r$、$\psi=\pm\dfrac{q}{2\pi}\theta$（正号源、负号汇），教材式 7.26、7.27 为源汇不在原点时的表达式。
\end{gongshi}

\begin{gongshi}{点汇的 $\varphi$、$\psi$（源流的逆过程；教材式 7.20--7.25 取负号）}
\[ u_r=-\frac{Q}{2\pi r},\qquad \varphi=-\frac{Q}{2\pi}\ln r,\qquad \psi=-\frac{Q}{2\pi}\theta \]
\textbf{含义：}四周流体沿平面均匀汇集到 $o$ 点被吸收，$o$ 点称为\textbf{汇点}。
点汇流动就是点源流动的逆过程，各表达式形式相同、\textbf{只差一个符号}，可直接写出。\\
\textbf{适用条件：}同点源；$Q$ 仍取正值，负号出现在公式里而不是数值里。\\
\textbf{流线形状：}指向汇点的半辐射线。
\end{gongshi}
\textbf{常考题型：}选择/填空（"源流的速度与半径成\rule{1.4cm}{0.4pt}比"$\to$ 反比）；
计算中作为叠加元件（半体＝均匀流＋点源）。
\textbf{重要度 \xstarfull{4}}\quad\yiju{E4 2015 年试题计算题第 8 题直接使用点源；E2"公式套用"}\quad\nandu{易}

\begin{banfa}
源汇只用一个尺寸判断：\textbf{"$u_r\propto 1/r$，所以速度有量纲上的 $1/r$ 尾巴"}。
若题目给的 $\varphi$ 里出现 $\ln r$（或 $\psi$ 里出现 $\theta$），\textbf{必是源或汇}，符号看流向；
若给的 $\varphi$ 里出现 $\theta$（或 $\psi$ 里出现 $\ln r$），\textbf{必是环流}；
若出现 $\cos\theta/r$（或 $\sin\theta/r$），\textbf{必是偶极子}。这三条认式子法能省下大量推导时间。
\end{banfa}
\end{kaodian}

\begin{kaodian}{4}{偶极子}
\begin{gongshi}{偶极子的 $\varphi$、$\psi$ 与速度场（教材式 7.37--7.40）}
\[ \varphi=+\frac{M\cos\theta}{2\pi r}=+\frac{Mx}{2\pi(x^2+y^2)},\qquad\qquad
   \psi=-\frac{M\sin\theta}{2\pi r}=-\frac{My}{2\pi(x^2+y^2)} \]
其中 $M=2aq$（$q$ 为源强，$2a$ 为源汇间距）：令 $2a\to0$、$q\to\infty$ 而乘积 $M=2aq$ 保持有限。
教材原话（书 p129）：“\textbf{设源点与汇点沿 $x$ 轴无限接近，即令 $2a\to0$，同时设 $q$ 无限增大，这样就能保证乘积 $2aq$ 始终保持为一常数 $M$：$M=2aq$。这一极限状态下，源汇合成的平面流动称为偶极子流，$M$ 称偶极子强度（$M>0$）}”。\\
\textbf{含义：}\textbf{偶极子＝源与汇无限接近的极限}：等强度的一个点源和一个点汇，
源汇间距 $\dd s\to0$ 的同时强度无限增大，使乘积 $M=q\,\dd s$ 趋于有限常数。
$M$ 称为\textbf{偶极矩}，其方向规定为\textbf{由汇指向源}（教材在 $(-a,0)$ 处放源、$(+a,0)$ 处放汇，故 $M$ 的方向就是 $-x$ 方向，书 p129）。\\
\textbf{适用条件：}不可压缩平面势流；除原点（奇点）以外处处无旋。\\
\textbf{流线形状：}在原点相切于 $x$ 轴的一族\textbf{圆}（$\psi=c$ 给出圆族），
等势线是与之正交的另一族圆。\\
\textbf{教材原话（书 p130）：}等势线“是圆心在 $y$ 轴上、\textbf{与 $y$ 轴相切于原点的圆}，$C>0$ 时圆位于 $y$ 轴右侧，否则在 $y$ 轴左侧”；流线“是圆心在 $\left(0,-\dfrac{M}{4\pi D}\right)$ 的圆，\textbf{圆与 $x$ 轴相切于原点}，当 $D>0$ 时圆位于 $x$ 轴下方，否则位于 $x$ 轴上方”。
\end{gongshi}
\textbf{常考题型：}填空（"偶极子由\rule{1.8cm}{0.4pt}的极限得到"$\to$ 等强度源汇无限接近）；
叠加题元件（均匀流＋偶极子＝圆柱体绕流）。
\textbf{重要度 \xstarfull{4}}\quad\yiju{E2"掌握基本概念＋公式套用"；E6 赵琴教材 7.2 节为独立小节（均匀流＋偶极子＝圆柱绕流，7.3 节直接调用）}\quad\nandu{中}

\begin{yicuo}
偶极矩方向极易记混：规定为\textbf{由汇指向源}（不是由源指向汇）。
另外记住：$\varphi=+\dfrac{M\cos\theta}{2\pi r}$、$\psi=-\dfrac{M\sin\theta}{2\pi r}$——\textbf{只有 $\psi$ 带负号，$\varphi$ 是正的}（教材式 7.39、7.40，书 p130）；
在"均匀流＋偶极子＝圆柱"的推导中，这个负号恰好用来把 $\psi$ 写成
$\psi=Ur\sin\theta\bigl(1-a^2/r^2\bigr)$ 的形状，所以\textbf{不能丢}；而分母里的 $2\pi$ 一旦丢掉，$a=\sqrt{M/(2\pi U)}$ 就再也对不上了。
\end{yicuo}

\begin{tishi}
\textbf{（对旧稿的订正·偶极子的 $\varphi$、$\psi$）}旧稿写 $\varphi=-\dfrac{M\cos\theta}{r}$、$\psi=-\dfrac{M\sin\theta}{r}$（$M=q\,\dd s$），有两处问题：\\
① \textbf{漏了分母中的 $2\pi$}，与本书“均匀流＋偶极子＝圆柱体绕流”一节所用的 $a=\sqrt{M/(2\pi U)}$ 自相矛盾；\\
② 这一对 $\varphi$、$\psi$ \textbf{不满足柯西-黎曼条件}（代入 $\partial\varphi/\partial x=\partial\psi/\partial y$ 不成立），而教材式 7.37 与 7.38 给出的一对 $\varphi=+\dfrac{Mx}{2\pi(x^2+y^2)}$、$\psi=-\dfrac{My}{2\pi(x^2+y^2)}$ 严格满足。\\
现统一按\textbf{教材式(7.37)--(7.40)（书 p129--130）}：源在 $(-a,0)$、汇在 $(+a,0)$，$M=2aq$，\textbf{$M$ 的方向由汇指向源}（即 $-x$ 方向、与来流相反，正好得到圆柱绕流）。\\
另说明（OCR 判读）：教材式 7.38、7.40 的\textbf{负号在扫描 OCR 中丢失}，但由式 7.37 的柯西-黎曼共轭关系、以及教材 7.3.3 节的推导过程都可判定为负号，故此处按负号采用。
\end{tishi}
\end{kaodian}

\begin{kaodian}{4}{纯环流（势涡）}
\begin{gongshi}{纯环流的 $\varphi$、$\psi$ 与速度场（教材式 7.30--7.32）}
\[ u_r=0,\qquad u_\theta=\frac{\Gamma}{2\pi r},\qquad
   \varphi=\frac{\Gamma}{2\pi}\theta,\qquad \psi=-\frac{\Gamma}{2\pi}\ln r \]
\textbf{含义：}流体绕某一点做匀速圆周运动、且速度与圆周半径\textbf{成反比}的流动。
衡量旋涡强度的物理量称为\textbf{速度环量} $\Gamma$，它不随圆周半径而变，具有方向性
（\textbf{逆时针为正}）。在平面势流中，环流不产生径向流动，只有切向速度。\\
\textbf{适用条件：}除原点外\textbf{处处无旋}——这一点是本章最反直觉的地方。
原点处 $u_\theta\to\infty$ 是唯一的有旋点（涡心），除它之外 $\bomega=0$，
所以纯环流是\textbf{有势、无旋}的流动，$\varphi$ 依然存在。\\
\textbf{流线形状：}以原点为圆心的\textbf{同心圆}；等势线为由原点出发的辐射线。
\end{gongshi}

\begin{gongshi}{环量与圆周速度的关系（教材式 7.30、7.83）}
\[ \Gamma=\oint_L \bu\cdot\dd\bl=\oint_L u_\theta r\,\dd\theta=2\pi r u_\theta
   \qquad\Longrightarrow\qquad u_\theta=\frac{\Gamma}{2\pi r} \]
\textbf{含义：}沿半径为 $r$ 的整个圆周积分，得 $\Gamma=2\pi r u_\theta$；
由它反解出速度场 $u_\theta=\Gamma/(2\pi r)$，这就是纯环流的来源。
$\Gamma$ 与 $r$ 无关，说明"环量"是刻画旋涡整体的量，而不是某处的速度。\\
\textbf{适用条件：}不可压缩流体（$\rho$ 可提出积分号）；$\Gamma$ 的正负按逆时针为正的约定。\\
\textbf{教材原话（书 p127）：}“\textbf{平面上流体质点绕一固定点作匀速圆周运动，不同半径圆周上质点运动速度反比于圆周半径，这就形成了平面上一点涡}”；“$\Gamma$ 在任一圆周上是一常数，称为点涡的强度，当 $\Gamma>0$ 时，表示质点作逆时针转动”。教材式 7.30--7.32 即 $u_r=0$、$u_\theta=\dfrac{\Gamma}{2\pi r}$、$\varphi=\dfrac{\Gamma}{2\pi}\theta$、$\psi=-\dfrac{\Gamma}{2\pi}\ln r$；教材式 7.35 给出复势 $W=\dfrac{\Gamma}{2\pi\mathrm{i}}\ln z$；教材式 7.33、7.34、7.36 分别为直角坐标表达式与点涡不在原点时的表达式。
\end{gongshi}

\begin{yicuo}
\textbf{（本章第三号易错点）把纯环流判成"有旋流动"。}\\
纯环流的流线是圆，但\textbf{流体微团本身并不旋转}：
取一个微元，外侧速度小、内侧速度大，两侧速度差刚好抵消了微团的转动，
算出 $w_z=\frac{1}{2}\left(\frac{\partial u_\theta}{\partial r}+\frac{u_\theta}{r}\right)-\frac{1}{2r}\frac{\partial u_r}{\partial\theta}
=\frac12\left(-\frac{\Gamma}{2\pi r^2}+\frac{\Gamma}{2\pi r^2}\right)=0$。\\
因此\textbf{"有环量"不等于"有旋"}：$\Gamma\ne0$ 而处处 $\bomega=0$，
纯环流恰恰是标准的有势无旋流（故又名\textbf{势涡}）。
有旋流动的典型是"刚体式旋转"（$u_\theta=\omega r$），它的微团真的在转。
\end{yicuo}

\begin{jielun}
四种基本平面势流一览（\textbf{背下这张表，本章大题就有了全部素材}）：
\begin{center}
\begin{tabularx}{\textwidth}{@{}lXXl@{}}
\toprule
流动 & 速度势 $\varphi$ & 流函数 $\psi$ & 流线形状 \\
\midrule
均匀流（沿 $x$，$u_x=U$） & $Ux$ & $Uy$ & 平行 $x$ 轴的直线 \\
点源（强度 $Q$） & $\dfrac{Q}{2\pi}\ln r$ & $\dfrac{Q}{2\pi}\theta$ & 由源点出发的辐射线 \\
点汇 & $-\dfrac{Q}{2\pi}\ln r$ & $-\dfrac{Q}{2\pi}\theta$ & 指向汇点的辐射线 \\
偶极子（偶极矩 $M=2aq$） & $+\dfrac{M\cos\theta}{2\pi r}$ & $-\dfrac{M\sin\theta}{2\pi r}$ & 与 $x$ 轴相切的圆族 \\
纯环流（环量 $\Gamma$） & $\dfrac{\Gamma}{2\pi}\theta$ & $-\dfrac{\Gamma}{2\pi}\ln r$ & 同心圆 \\
\bottomrule
\end{tabularx}
\end{center}
\textbf{记忆规律：}\textbf{$\ln r$ 与 $\theta$ 是一对（源汇、环流）}，
\textbf{$\cos\theta/r$ 与 $\sin\theta/r$ 是一对（偶极子）}，
\textbf{均匀流是"没有 $r$ 的线性项"}。凡 $\varphi$ 中出现 $\ln r$、$\varphi$ 中出现 $\theta$、
$\psi$ 中出现 $\theta$、$\psi$ 中出现 $\ln r$ 的组合方式，一律按上表对照。
\end{jielun}

\textbf{常考题型：}填空（给出某势流的 $\varphi$ 或 $\psi$，问"这是哪一种基本流动"）；
选择（"点源流的速度与半径成\rule{1.2cm}{0.4pt}比"）；简答（"什么是势涡，它是有旋还是无旋"）。
\textbf{重要度 \xstarfull{4}}\quad\yiju{E4 2015 年试题计算题第 8 题使用了点源；E2"第六至第八章掌握基本概念＋公式套用即可"；E6 教材 7.2 节}\quad\nandu{易}
\end{kaodian}

\section{势流的叠加}

\begin{kaodian}{4}{势流叠加原理}
\begin{gongshi}{叠加原理（教材 7.3.1 节，书 p129，教材未给式号）}
\[ \varphi=\varphi_1+\varphi_2+\cdots+\varphi_n,\qquad\qquad \psi=\psi_1+\psi_2+\cdots+\psi_n \]
\textbf{含义：}把几个简单势流叠加成一个复合流动时，复合流动的流函数（或势函数）
就是各原始流动的流函数（或势函数）的\textbf{简单代数相加}。
之所以成立，是因为拉普拉斯方程是\textbf{线性}的：若 $\grad^2\varphi_1=0$、$\grad^2\varphi_2=0$，
则 $\grad^2(\varphi_1+\varphi_2)=0$——两个调和函数的和仍是调和函数。\\
\textbf{适用条件：}各叠加流动必须是\textbf{同一不可压缩流体的无旋流动}，且叠加后要检查叠加是否破坏了
原有的无旋性（只在与原奇点重合处例外）。叠加后速度场也是各速度场的矢量和。\\
\textbf{教材原话（书 p129）：}“设想有 $n$ 个简单平面势流……现将 $n$ 个平面流动叠加得到一个新的平面流动，\textbf{新的流动仍然是一有势流动}”，其势函数为 $\varphi=\varphi_1+\varphi_2+\cdots+\varphi_n$、流函数为 $\psi=\psi_1+\psi_2+\cdots+\psi_n$；教材还给出\textbf{复势形式}：“它们的和 $W=W_1+W_2+\cdots+W_n$ 依然为一解析的复变函数，仍可作为某一有势流动的复势”；并指出得到 $\varphi$、$\psi$ 后“可求出流场的速度矢量，进一步以伯努利方程求出压强分布，完成流场分析”。
\end{gongshi}
\textbf{常考题型：}填空（"势流叠加原理的依据是控制方程是\rule{1.6cm}{0.4pt}方程"$\to$ 线性／拉普拉斯）；
计算（叠加后求驻点、求流线，见下）。
\textbf{重要度 \xstarfull{4}}\quad\yiju{E4 2015 年试题计算题第 8 题（均匀流叠加点源）；E2"公式套用"；E6 教材 7.3 节}\quad\nandu{中}
\end{kaodian}

\begin{kaodian}{2}{教材 7.3.2 节：点源（汇）＋点涡＝螺旋流}
\textbf{为什么补这一节：}教材在 7.3.2 节用一整节讲“点源（汇）与点涡—螺旋流”，并给出工程实例（水泵蜗壳）；
而教材\textbf{没有}“均匀流＋点源＝半体”这样的独立小节（按已 OCR 的本章 31 页正文检索，无“半体”“兰肯”字样）。
两块都要会：\textbf{教材考螺旋流，真题考半体}。

\begin{gongshi}{螺旋流：点汇＋点涡的叠加（教材 7.3.2 节，书 p129，教材未给式号）}
\[ u_r=\frac{q}{2\pi r},\qquad u_\theta=\frac{\Gamma}{2\pi r},\qquad
   \frac{u_\theta}{u_r}=\frac{\Gamma}{q}=\text{常数},\qquad u=\frac{\sqrt{q^2+\Gamma^2}}{2\pi r} \]
\[ \varphi=\frac{q}{2\pi}\ln r+\frac{\Gamma}{2\pi}\theta,\qquad\qquad
   \psi=\frac{q}{2\pi}\theta-\frac{\Gamma}{2\pi}\ln r \]
\textbf{含义：}把一个强度为 $q$ 的点汇（或点源）与一个强度为 $\Gamma$ 的点涡叠加：流体质点一面被“吸”向中心、一面绕中心旋转，
合速度 $u$ 与径向的夹角 $\alpha$ 保持常数（$\tan\alpha=\Gamma/q$），而 $u$ 仍与 $r$ 成反比。\\
\textbf{教材原话（书 p129）：}“平面坐标原点有一强度为 $q$ 的点汇和一强度为 $\Gamma$ 的点涡（$q>0$，$\Gamma>0$）”；
“令 $\psi=C$，经计算可以得到 $r=C_1\mathrm{e}^{q\theta/\Gamma}$，在平面极坐标系下，这一方程代表的\textbf{等势线是平面对数螺旋线}，同样可以得到\textbf{流线也是平面对数螺旋线}，这两条曲线是正交的”；“\textbf{水泵矩形断面蜗壳中的理想流动是点源与点涡叠加的一个例子}”。\\
\textbf{适用条件：}不可压缩平面势流；点汇（源）与点涡同在坐标原点；$q>0$、$\Gamma>0$（逆时针）。\\
\textbf{OCR 说明：}教材 7.3.2 节的 $\varphi$、$\psi$ 公式在扫描件中只读到“$\ln r+$”“$\theta$、$\ln r$”等残形且未编号，
上式按 7.3.1 节的叠加原理与柯西-黎曼条件写出，符号取教材“点汇＋逆时针点涡”。
\end{gongshi}

\textbf{常考题型：}选择/填空（“点源（汇）与点涡叠加得到\rule{1.8cm}{0.4pt}”$\to$ 螺旋流；“流线是\rule{1.6cm}{0.4pt}”$\to$ 对数螺旋线）；
简答（“水泵蜗壳内的理想流动是哪两种基本势流的叠加”$\to$ 点源（汇）与点涡）。
\textbf{重要度 \xstarfull{2}}\quad\yiju{E2“第六到第八章只有少量填空/判断/选择题，掌握基本概念＋公式套用即可”；E6 教材 7.3.2 节（书 p129）教材为独立小节且有工程实例}\quad\nandu{中}
\end{kaodian}

\begin{tishi}
\textbf{（教材与旧稿的对照）}旧稿写的是“均匀流＋点源＝兰肯半体”，教材 7.3 节讲的是
① 叠加原理（7.3.1）、② 点源（汇）＋点涡＝\textbf{螺旋流}（7.3.2）、③ 偶极子流（7.3.3），\textbf{没有兰肯半体}。
但 2015 年计算题第 8 题考的正是“均匀流叠加点源求驻点与流线”，所以下面“兰肯半体”一框保留（依据是 E4 真题），
其公式按叠加原理自行推得；考试时\textbf{教材的螺旋流与真题的半体都要会}。
\end{tishi}

\begin{kaodian}{4}{驻点与流线的求法，及兰肯半体}
\begin{gongshi}{驻点的定义与求法}
\[ \text{驻点：}\quad u_x=0\ \text{且}\ u_y=0\qquad\Longleftrightarrow\qquad \grad\varphi=0 \]
\textbf{含义：}速度为零的点称为\textbf{驻点}（停滞点）。在 $\varphi$、$\psi$ 已知时，
对速度分量列两个方程组即可解出驻点坐标；用极坐标时把 $u_r$、$u_\theta$ 同时置零更省事。\\
\textbf{适用条件：}驻点必须落在\textbf{实际流场中}（物理上可达的区域）；
解出的点若位于源点等奇点上或流场外，要舍去。
\end{gongshi}

\begin{gongshi}{兰肯半体：均匀流 ＋ 点源}
\[ \varphi=Ux+\frac{Q}{2\pi}\ln r,\qquad\qquad
   \psi=Uy+\frac{Q}{2\pi}\theta=U r\sin\theta+\frac{Q}{2\pi}\theta \]
\[ \text{驻点：}\quad x_s=-\frac{Q}{2\pi U}=-a,\qquad\qquad
   \text{半体下游无穷远处的半宽度：}\quad y_\infty=\frac{Q}{2U} \]
\textbf{含义：}在沿 $+x$ 的均匀流中于原点放置一个强度为 $Q$ 的点源，两者叠加即\textbf{兰肯半体}。
沿 $x$ 轴置 $u_x=U+\dfrac{Q}{2\pi x}=0$ 得 $x_s=-\dfrac{Q}{2\pi U}$，
即驻点在\textbf{源点上游}距离 $a=\dfrac{Q}{2\pi U}$ 处。
过驻点的那条流线（代入 $x=-a$、$y=0$、$\theta=\pi$ 得 $\psi_s=Q/2$）的上半支即为\textbf{半体轮廓}：
它把流场分成两部分——$\psi\ge Q/2$ 的流体从无穷远流来、绕过这个"畸变体"再回到无穷远，
$\psi<Q/2$ 的部分被"困"在内部，相当于一个无限长的半椭球状固体障碍物，
所以"均匀流＋点源"是用基本势流"造"出固体壁面模型的经典例子。
在 $\theta\to0$、$r\to\infty$ 处由 $Uy+Q\theta/(2\pi)=\psi_s$ 得 $y_\infty=\psi_s/U=\dfrac{Q}{2U}$。\\
\textbf{适用条件：}不可压缩平面势流；点源位于原点、来流沿 $+x$ 方向、速度大小为 $U$。
\end{gongshi}

\begin{banfa}
\textbf{叠加题统一四步（2015 年第 8 题的标准解法）：}
\begin{enumerate}
  \item \textbf{写出 $\varphi$、$\psi$：}把各基本流动的 $\varphi$、$\psi$ 代数相加，
        尽量换成直角坐标与 $r$、$\theta$ 混用（求驻点用直角，求形状用极坐标）。
  \item \textbf{求驻点：}令 $u_x=u_y=0$ 解方程组；
        "均匀流＋点源"只需令 $u_x=0$ 即可（因 $u_y=Qy/2\pi r^2$ 在 $y=0$ 上为零），
        得 $x_s=-Q/(2\pi U)$。
  \item \textbf{求过驻点的流线：}把驻点坐标代回 $\psi$，得到 $\psi_s$，
        写"过驻点的流线方程为 $\psi=\psi_s$"即可，\textbf{不必解出显式 $y(x)$}。
  \item \textbf{回答"下游无穷远半宽度"（如题目追问）：}在 $\theta\to0$、$r\to\infty$ 处用
        $Uy+\frac{Q}{2\pi}\theta=\psi_s$，令 $\theta\to0$ 得 $y_\infty=\psi_s/U=Q/(2U)$。
\end{enumerate}
\textbf{识别技巧：}"求驻点位置及经过驻点的流线方程"这类问法的答案永远是两步——
一个坐标＋一个 $\psi=$ 常数的方程，不要把 $\psi$ 展开成 $y=y(x)$ 的复杂形式再解。
\end{banfa}

\textbf{常考题型：}计算（2015 年计算题第 8 题原文如下）。
\textbf{重要度 \xstarfull{4}}\quad\yiju{E4 2015 年试题计算题第 8 题；E6 教材 7.3.1 节叠加原理（教材无“半体”专节，本条公式按叠加原理推得）}\quad\nandu{中}

\begin{zhenti}{2015 年试题 计算题第 8 题}{计算题（12 分）}
将速度为 $v_\infty$ 的平行于 $x$ 轴的均匀流和在原点强度为 $q$ 的点源叠加，
求叠加后流场中驻点位置及经过驻点的流线方程。
\end{zhenti}
\end{kaodian}

\begin{kaodian}{3}{均匀流 ＋ 偶极子 ＝ 圆柱体绕流}
\begin{gongshi}{叠加结果与驻点（教材式 7.41--7.45）}
\[ \psi=U r\sin\theta\left(1-\frac{a^2}{r^2}\right),\qquad\text{其中圆半径 } a=\sqrt{\frac{M}{2\pi U}} \]
\textbf{含义：}沿 $+x$ 的均匀流 $U$ 叠加一个偶极矩方向与来流相反（偶极子指向 $-x$、轴与 $x$ 轴重合）
的偶极子，取 $a^2=M/(2\pi U)$，则 $\psi=0$ 的流线由两部分组成：
\textbf{$r=a$ 的圆}（即圆柱表面）与 $x$ 轴。
$r=a$ 这条流线把流场分成"圆外"和"圆内"两部分，圆内解无物理意义、
圆外的解就代表\smash{均匀流绕半径为 $a$ 的圆柱体的绕流}。
\textbf{驻点}为 $r=a$、$\theta=0$ 与 $\theta=\pi$ 两点（即 $(\pm a,0)$），恰在圆柱表面的前后两点。\\
\textbf{适用条件：}不可压缩平面势流；偶极矩方向必须与来流方向相反（若取同向，得到的仍是同样的圆，
但流向要按速度场重新判定）。\\
\textbf{教材原话（书 p130--131）：}“\textbf{设想将圆柱体从流场中抽去，然后在坐标原点处设置一强度 $M=2\pi u_\infty r_0^2$ 的偶极子}”（教材式 7.43、7.44 即 $\varphi=u_\infty r\cos\theta\left(1+\dfrac{r_0^2}{r^2}\right)$、$\psi=u_\infty r\sin\theta\left(1-\dfrac{r_0^2}{r^2}\right)$，教材式 7.45 为复势 $W=u_\infty z+\dfrac{M}{2\pi z}$）；
教材指出 $\psi=0$ 的两支为 $y=0$ 与 $x^2+y^2=r_0^2$，并强调“\textbf{流体不可能穿越流线，因而理想的流线可以与固体壁面互换}”（书 p131）。
\end{gongshi}
\textbf{常考题型：}选择/填空（"均匀流与偶极子叠加得到\rule{1.8cm}{0.4pt}绕流"$\to$ 圆柱体）；
与下一节（圆柱体绕流的压强与阻力）连成一道完整大题。
\textbf{重要度 \xstarfull{3}}\quad\yiju{E2"第六至第八章只有少量填空/判断/选择题，掌握基本概念＋公式套用即可"；E1 slide33"概念为主"}\quad\nandu{中}
\end{kaodian}

\begin{kaodian}{3}{点汇 ＋ 点源 ＋ 偶极子 ＝ 圆柱体带环流绕流}
\begin{gongshi}{含环流的复合流动（教材式 7.51--7.58）}
\[ \psi=U r\sin\theta\left(1-\frac{a^2}{r^2}\right)-\frac{\Gamma}{2\pi}\ln r \]
\textbf{含义：}在"均匀流＋偶极子"（即圆柱绕流）的 $\psi$ 上再叠加一个纯环流
（环量为 $\Gamma$），得到有环量的圆柱绕流。$\Gamma=0$ 时退化为上一节的无环量圆柱绕流。
具体闭环流动的产生方式（如绕圆柱旋转）不在考试范围，只需知道\textbf{"环流叠加改变流线、产生升力"}。\\
\textbf{适用条件：}不可压缩平面势流；环流为逆时针时 $\Gamma>0$。\\
\textbf{（对旧稿的订正·环流项符号）}旧稿写 $\psi=Ur\sin\theta\left(1-\dfrac{a^2}{r^2}\right)+\dfrac{\Gamma}{2\pi}\ln r$，\textbf{环流项的符号反了}：
教材式 7.32 纯环流的 $\psi=-\dfrac{\Gamma}{2\pi}\ln r$，教材式 7.54 圆柱有环量绕流的
$\psi=u_\infty r\sin\theta\left(1-\dfrac{r_0^2}{r^2}\right)-\dfrac{\Gamma}{2\pi}\ln r$（书 p133），两处都是\textbf{负号}，故已改为减号。\\
\textbf{教材式号（书 p133）：}有环量绕流的 $\varphi$、$\psi$ 直角坐标式为教材式 7.51、7.52，极坐标式为教材式 7.53、7.54，复势为教材式 7.55（$W=u_\infty z+\dfrac{M}{2\pi z}+\dfrac{\Gamma}{2\pi\mathrm{i}}\ln z$），速度分布为教材式 7.56、7.57，柱面上速度为教材式 7.58（$u_r=0$、$u_\theta=-2u_\infty\sin\theta+\dfrac{\Gamma}{2\pi r_0}$）。
\end{gongshi}
\textbf{常考题型：}填空/选择（叠加结果的组成）；与库塔-儒可夫斯基定理合并考"升力来自什么"。
\textbf{重要度 \xstarfull{3}}\quad\yiju{E2"掌握基本概念＋公式套用"；E6 教材 7.3、7.4 节}\quad\nandu{中}
\end{kaodian}

\section{圆柱体绕流}

\subsection{无环流圆柱绕流与达朗贝尔佯谬}

\begin{kaodian}{3}{无环量圆柱绕流：$\varphi$、$\psi$、表面速度与压强分布}
\begin{gongshi}{无环量圆柱绕流的 $\varphi$、$\psi$ 与速度场（教材式 7.41--7.47）}
\[ \varphi=U\left(r+\frac{a^2}{r}\right)\cos\theta,\qquad\qquad
   \psi=U\left(r-\frac{a^2}{r}\right)\sin\theta \]
\[ u_r=U\left(1-\frac{a^2}{r^2}\right)\cos\theta,\qquad
   u_\theta=-U\left(1+\frac{a^2}{r^2}\right)\sin\theta \]
\textbf{含义：}$U$ 为来流速度、$a$ 为圆柱半径、$r\ge a$ 为场点半径。
圆柱表面 $r=a$ 处 $u_r=0$（流体不能穿过壁面，符合不穿透边界条件），
只剩切向速度 $u_\theta=-2U\sin\theta$。\\
\textbf{适用条件：}不可压缩理想流体、平面无环量绕流、$r\ge a$。
$\theta$ 由 $+x$ 轴（下游方向）量起；$\theta=0,\ \pi$ 为驻点（前后停驻点）。
\end{gongshi}

\begin{gongshi}{圆柱表面的速度最大值与压强分布（教材式 7.46--7.50）}
\[ |u_{\theta}|_{\max}=2U\ \ \text{（在 }\theta=\pm90^\circ\text{ 处，即圆柱上下两侧顶点）} \]
\[ p=p_\infty+\frac12\rho U^2\left(1-4\sin^2\theta\right) \]
\textbf{含义：}由伯努利方程（无黏、不可压缩、沿流线）在无穷远与圆柱表面之间列式：
$p_\infty+\frac12\rho U^2=p+\frac12\rho u_\theta^2$，
代入 $u_\theta=-2U\sin\theta$ 即得表面压强分布式。
$\theta=\pm90^\circ$ 处 $u_\theta$ 最大（$2U$）、压强最低；
$\theta=0,\pi$（前、后驻点）处 $u_\theta=0$、压强最高，为 $p_\infty+\frac12\rho U^2$；
压强分布关于 $x$ 轴与 $y$ 轴均对称。\\
\textbf{适用条件：}理想（无黏）不可压缩流体、无环量、圆柱表面上的点；
用伯努利方程的前提是忽略质量力（或沿水平流线）且流动无旋。\\
\textbf{教材原话（书 p131--132）：}“在 $\theta=0$（B 点）和 $\theta=2\pi$（A 点）处 $u_\theta=0$，A、B 两点是\textbf{分流点，也称为驻点}。在 $\theta=\pm90^\circ$ 处 $u_\theta$ 达到最大值 $|u_\theta|=2u_\infty$”；“在圆柱表面的两个驻点处，压强达到最大值，而在圆柱表面 $\theta=\pm90^\circ$ 处压强降到最低”；“\textbf{沿圆柱面的压强系数与圆柱体的半径以及均匀流的速度、压强分布无关}”。教材式 7.49、7.50 即 $C_p=\dfrac{p-p_\infty}{\frac12\rho u_\infty^2}=1-4\sin^2\theta$。
\end{gongshi}

\textbf{常考题型：}填空（"圆柱表面最大速度为来流的\rule{1.2cm}{0.4pt}倍"$\to$ 2 倍）；
计算（给出 $U$、$a$、$p_\infty$，求某点表面压强或驻点压强差）；
简答（"为什么圆柱绕流阻力为零"，见下）。
\textbf{重要度 \xstarfull{3}}\quad\yiju{E1 slide33"第七章＝概念为主"；E2"第六至第八章只有少量填空/判断/选择题"；E6 教材 7.4 节}\quad\nandu{中}

\begin{banfa}
圆柱绕流的题几乎都围绕三个数转，\textbf{记三个"整倍数"}：
\textbf{最大速度＝2$U$}（在 $\pm90^\circ$）；\textbf{驻点压强比远场高 $\frac12\rho U^2$}；
\textbf{最低压强比远场低 $\frac32\rho U^2$}（$\theta=\pm90^\circ$ 处，把 $\sin^2\theta=1$ 代入
$p=p_\infty+\frac12\rho U^2(1-4)=p_\infty-\frac32\rho U^2$）。
题目给"求最低压强出现的位置及数值"，答"$\theta=\pm90^\circ$，$p_\infty-\frac32\rho U^2$"即可，
不必重推伯努利方程。
\end{banfa}
\end{kaodian}

\begin{kaodian}{3}{达朗贝尔佯谬（阻力为零）与库塔-儒可夫斯基升力}
\begin{gongshi}{达朗贝尔佯谬：无环量时阻力为零（教材 7.4.1 节柱面合力，书 p132）}
\[ F_D=\int_0^{2\pi}-p\cos\theta\cdot a\,\dd\theta=0 \]
\textbf{含义：}作用在圆柱上的合力由表面压强积分得到。
$p=p_\infty+\frac12\rho U^2(1-4\sin^2\theta)$ 是 $\theta$ 的偶函数、且大小相等地
分布在前后（$\theta$ 与 $\pi-\theta$）对应点上，于是 $x$ 方向（来流方向）的压强合力恒为零：
\textbf{无环量时圆柱在理想流体中不受阻力}，这一反直觉结论称为\textbf{达朗贝尔佯谬}。\\
\textbf{适用条件：}\textbf{理想（无黏）流体}、无环量、定常、不可压缩。
真实流体中因黏性引起\textbf{边界层分离与尾涡}，圆柱前后压强不再对称，阻力并不为零。\\
\textbf{教材原话（书 p132）：}“以 $\theta_0$ 代替 $\theta$，压强值 $p$ 不变，说明圆柱体表面压强分布对称于 $x$ 轴……以 $\pi-\theta_0$ 代替 $\theta$，压强值也不变，说明圆柱体表面压强分布对称于 $y$ 轴……\textbf{这样，圆柱表面流体压强的合力为零}”；“理想均匀流绕流圆柱体时，圆柱体的\textbf{阻力 $D$ 和升力 $L$ 都等于零}”。
\end{gongshi}

\begin{gongshi}{库塔-儒可夫斯基升力定理（教材式 7.61，书 p135）}
\[ L=\rho U\Gamma \]
\textbf{含义：}在均匀流叠加环量 $\Gamma$ 后，压强分布不再关于 $x$ 轴对称，
垂直于来流方向出现合力——\textbf{单位长度圆柱所受的升力}为 $L=\rho U\Gamma$，
方向由"来流方向旋转 $90^\circ$ 且与环量方向相反"定出：
\textbf{来流沿 $+x$、环量逆时针（$\Gamma>0$）时，升力向下}（沿 $-y$）。\\
\textbf{教材原话（书 p135）：}“上式称为库塔-儒可夫斯基升力公式，式中\textbf{负号表明圆柱体所受流体升力方向沿 $y$ 轴负向}”；“升力大小与来流速度 $u_\infty$、流体密度 $\rho$ 和旋转圆柱引起的环量 $\Gamma$ 成正比，\textbf{升力方向应将来流方向沿圆柱旋转方向反向旋转 $90^\circ$}”；该公式“可以推广应用于理想流体均匀流绕过任何形状有环量无分离的平面流动，例如具有流线型的翼型绕流等，在轴流式水泵、水轮机叶片设计中有重要应用”。\\
\textbf{（对旧稿的订正·升力方向）}旧稿此处写“$\Gamma>0$ 时升力向上（沿 $+y$）”，\textbf{方向反了}。由教材式 7.58（$u_\theta=-2u_\infty\sin\theta+\dfrac{\Gamma}{2\pi r_0}$）可见 $\Gamma>0$ 时\textbf{圆柱上部速度减小、下部速度增大}，据伯努利方程上部压强高于下部，合力指向 $-y$——这正是教材式 7.61 出现负号的原因，也与“来流方向沿旋转方向反向旋转 $90^\circ$”一致（$+x$ 反向转 $90^\circ$ 即 $-y$）。\\
\textbf{适用条件：}理想不可压缩流体的定常平面势流，含环量的圆柱（或任意截面柱体）绕流；
$L$ 为\textbf{单位长度}（单位展长）上的升力，方向与来流垂直。\\
\textbf{驻点位置随 $\Gamma$ 变化（教材式 7.59，书 p134）：}$\sin\theta=\dfrac{\Gamma}{4\pi u_\infty r_0}$，$\theta$ 由 $+x$ 轴逆时针量起。\\
\textbf{（对旧稿的订正·驻点移动方向）}旧稿写“\textbf{逆时针}环量使驻点\textbf{下移}”，\textbf{与教材正好相反}。教材原话（书 p134）：“当柱体作\textbf{顺时针}转动，点涡强度 $\Gamma<0$ 时……使驻点\textbf{向下}移动”；“\textbf{当柱体作逆时针转动，驻点的位置沿 $y$ 轴向上移动}”。\\
\textbf{教材给出的三种情形（书 p134）：}① $|\Gamma|<4\pi u_\infty r_0$：两个驻点仍在圆柱表面上并对称于 $y$ 轴，随 $|\Gamma|$ 增大\textbf{向下移动并互相靠拢}；② $|\Gamma|=4\pi u_\infty r_0$：两个驻点\textbf{重合成一点，出现在圆柱表面的最下端}；③ $|\Gamma|>4\pi u_\infty r_0$：柱面上不再有驻点，\textbf{驻点脱离柱面沿 $y$ 轴下移到某一位置}。\\
\textbf{考场判断法：}直接把 $\Gamma$ 的正负代进 $\sin\theta=\dfrac{\Gamma}{4\pi u_\infty r_0}$——$\Gamma>0$ 时 $\sin\theta>0$，驻点在\textbf{上半柱面}（上移）；$\Gamma<0$ 时在下半柱面（下移）。教材讲完 $\Gamma<0$ 之后写的“随 $|\Gamma|$ 值的增加而向下移动，并互相靠拢”，是以 $\Gamma<0$ 为前提的，考试以公式符号为准。
\end{gongshi}

\begin{yicuo}
\textbf{（本章第四号易错点）把达朗贝尔佯谬的结论无限推广。}\\
"理想流体绕圆柱无环流时阻力为零"是对的，但\textbf{不能推出}"理想流体绕任何物体阻力都为零"
或"实际流体绕圆柱阻力也为零"。达朗贝尔佯谬与实验不符，\textbf{根本原因是忽略了黏性}——
真实流体在圆柱表面形成边界层，逆压区发生\textbf{边界层分离}、
下游形成相反旋转的\textbf{尾涡}，前后压强分布不再对称，于是产生了压差阻力（形状阻力）。
另外：\textbf{升力与阻力不是一回事}——升力靠\textbf{环量}产生（库塔-儒可夫斯基），
阻力在理想流体中为零、在真实流体中由黏性产生。
\end{yicuo}

\textbf{常考题型：}简答（"为什么理想流体绕圆柱无环量时阻力为零？该结论为何与实验不符？"）；
填空（"无环量圆柱绕流不产生\rule{1.4cm}{0.4pt}"$\to$ 阻力；"库塔-儒可夫斯基升力 $L=$ \rule{1.6cm}{0.4pt}"$\to$ $\rho U\Gamma$）；
选择（升力方向与来流、环量方向的关系）。
\textbf{重要度 \xstarfull{3}}\quad\yiju{E1 slide33"第七章＝概念为主"；E2"第六到第八章只有少量填空/判断/选择，掌握基本概念即可"；E6 教材 7.4 节}\quad\nandu{中}
\end{kaodian}

\section{空间势流}

\begin{kaodian}{1}{空间势流简介：空间点源汇、偶极子、绕球流动与附加质量}
\begin{gongshi}{空间点源（三维源）的 $\varphi$ 与速度场（教材式 7.64、7.65）}
\[ u_r=\frac{Q}{4\pi r^2},\qquad\qquad \varphi=-\frac{Q}{4\pi r} \]
\textbf{含义：}流体由空间中一点沿半径方向均匀流出的流动。
与平面源的最大区别是\textbf{球面积 $\propto r^2$}，所以速度为 $1/r^2$ 衰减、
$\varphi$ 为 $-Q/(4\pi r)$（平面源是 $1/r$ 与 $\ln r$）。
$Q$ 为体积流量。\\
\textbf{适用条件：}不可压缩无旋空间流动；$r>0$，源点为奇点。\\
\textbf{教材原话（书 p137）：}教材用球坐标系（$R,\theta,\beta$）给出空间点源的速度分量 $u_R=\pm\dfrac{q}{4\pi R^2}$、$u_\theta=u_\beta=0$，积分得势函数（教材式 7.65）$\varphi=\mp\dfrac{q}{4\pi R}$；教材式 7.64 为空间均匀流 $\varphi=u_\infty R\cos\theta$。
此处教材的 $R$ 与本节（以及本章其余部分）的 $r$ 含义相同，都是到场点的距离。\\
\textbf{轴对称流函数（教材 7.5.3 节，书 p139--140，旧稿完全没写）：}教材用 7.5.2 节的轴对称流函数给出
空间均匀流 $\psi=\dfrac12 u_\infty R^2\sin^2\theta$（教材式 7.71）、空间点源 $\psi=-\dfrac{q}{4\pi}\cos\theta$（教材式 7.72）、
空间偶极子 $\psi=-\dfrac{M}{4\pi R}\sin^2\theta$（教材式 7.73）——这是本章流函数不能漏的一块（另见前面“轴对称流动有没有流函数”的订正框）。
\end{gongshi}

\begin{gongshi}{空间偶极子与绕球流动}
\[ \varphi_{\text{偶极子}}=+\frac{M\cos\theta}{4\pi r^2},\qquad\qquad
   \varphi_{\text{绕球}}=U\left(r+\frac{a^3}{2r^2}\right)\cos\theta \]
\textbf{含义：}空间偶极子＝三维源汇无限接近的极限，$\varphi$ 按 $1/r^2$ 衰减；
均匀流叠加适当强度的时空偶极子，得到\textbf{速度为 $U$ 的均匀流绕过半径 $a$ 的球}的流动，
球面上 $u_r=0$，最大速度出现在球的赤道处。\\
\textbf{适用条件：}不可压缩理想流体、无旋、定常；球面不穿透条件 $u_r\big|_{r=a}=0$。\\
\textbf{教材原话（书 p138）：}空间偶极子由“在空间流动中，等强度的点源和点汇叠加”取极限得到，“首先将空间点汇置于 $+z$ 轴上，点源置于 $-z$ 轴上”，“即满足 $\lim q\Delta l=M$，$M$ 为一常数值，这样就可以得到空间偶极子流，$M$ 称为空间偶极子的强度（或偶极矩）”，其势函数为教材式 7.66。\\
\textbf{（对旧稿的订正·空间偶极子）}旧稿写 $\varphi_{\text{偶极子}}=-\dfrac{M\cos\theta}{r^2}$，与教材式 7.66 相比\textbf{漏了 $4\pi$、且符号相反}；现按教材式 7.66 改为 $+\dfrac{M\cos\theta}{4\pi r^2}$（教材取 $M=\lim q\Delta l$，与平面偶极子 $M=2aq$ 的定义不同）。\\
\textbf{圆球绕流的教材结论（教材式 7.74--7.78，书 p140--141）：}由零流面得 $\dfrac{M}{2\pi u_\infty}=a^3$，即 $M=2\pi u_\infty a^3$；
$\psi=u_\infty r^2\left[1-\left(\dfrac{a}{r}\right)^3\right]\sin^2\theta$、$\varphi=u_\infty r\left[1+\dfrac12\left(\dfrac{a}{r}\right)^3\right]\cos\theta$；
球面速度分布 $u_r=0$、$u_\theta=\dfrac32 u_\infty\sin\theta$，教材原话：“最大速度发生在 $\theta=\pm\pi/2$，
$|u_\theta|_{\max}=\dfrac32 u_\infty$”“\textbf{绕圆球表面的速度最大值不如绕圆柱大}，这是因为绕圆球时流体有较宽裕的空间流过物体”；
球面压强系数 $C_p=1-\dfrac94\sin^2\theta$，教材原话：“圆球表面上的压强分布是关于水平轴和铅直轴对称的，因而其合力等于零。因此，\textbf{圆球被均匀流绕过时不受流体合力作用}”。
\end{gongshi}

\noindent\textbf{附加质量（虚拟质量）：}物体在流体中做\textbf{非定常}运动时，必须带动周围流体一起加速，
流体因此对物体产生一个附加的惯性作用，其效果相当于物体的质量增加了若干——
这个等效增加的质量称为\textbf{附加质量（虚拟质量）}。
记为 $m_{\text{附}}=\rho V k$（$k$ 取决于物体形状，球体 $k=0.5$）。\\

\begin{tishi}
\textbf{（教材覆盖情况说明）}在已 OCR 的本章 31 页正文（书 p119--149）中\textbf{没有出现“附加质量”或“虚拟质量”}——教材 7.5 节讲的是空间点源汇、轴对称流函数、圆球绕流与回转体（轴对称体）绕流。
本段内容来自外校辅导书整理稿（杨树人《工程流体力学考点精讲》），真题中也无对应题目，故重要度保持 \xstarfull{1}；
时间紧时只需记住一句：\textbf{物体作非定常运动要带动周围流体一起加速，相当于物体质量增大了若干}。
\end{tishi}
\textbf{常考题型：}概念选择（"附加质量反映的是\rule{1.6cm}{0.4pt}"$\to$ 带动周围流体加速的惯性效应）。
\textbf{重要度 \xstarfull{1}}\quad\yiju{E4 历年真题中未见本章内容成题；E1 slide33"第七章＝概念为主"（仅章节级依据）}\quad\nandu{中}
\end{kaodian}

\section{理想流体的旋涡运动}

\begin{kaodian}{3}{涡线、涡管、涡束与涡通量}
\textbf{是什么：}流体做有旋流动时，其质点存在旋转，有时称为涡流。
描述有旋流动的几何与强度概念有四个。

\begin{center}
\begin{tabularx}{\textwidth}{@{}lXX@{}}
\toprule
概念 & 定义（教材口径） & 关键点 \\
\midrule
涡线 & 曲线上各质点在同一瞬间的\textbf{旋转角速度矢量与曲线在该点相切} & 与流线的定义方式完全平行，只是把速度换成角速度 \\
涡束（涡管） & 通过微元断面的涡线组成涡束，涡束的\textbf{表面}称为涡管 & 涡管是一根"管子"，管壁上角速度与管壁相切 \\
涡通量 $I$ & 涡束断面面积与 $2$ 倍旋转角速度的乘积，$\dd I=2w\,\dd A=\Omega\,\dd A$ & $\Omega=2w$ 称为\textbf{旋度}，也称\textbf{涡量} \\
速度环量 $\Gamma$ & $\Gamma=\oint_L\bu\cdot\dd\bl=\oint_L\left(u_x\dd x+u_y\dd y+u_z\dd z\right)$ & 逆时针为正；$\dd\bl$ 为沿封闭曲线 $L$ 的有向线元 \\
\bottomrule
\end{tabularx}
\end{center}

\begin{gongshi}{斯托克斯定理：速度环量与涡通量的关系（教材式 7.83、7.84）}
\[ \Gamma=\oint_L \bu\cdot\dd\bl=\iint_A \bomega\cdot\dd\bA=\iint_A \Omega\,\dd A \]
\textbf{含义：}沿封闭曲线 $L$ 的速度环量，等于穿过以 $L$ 为边界的曲面 $A$ 的\textbf{涡通量}。
流体力学里的斯托克斯定理是数学上格林公式／斯托克斯公式在速度场与旋度上的应用，
沟通了\textbf{线积分}（环量）与\textbf{面积分}（涡通量）两条路。\\
\textbf{适用条件：}速度场连续可微、曲面 $A$ 以 $L$ 为边界且完全落在流场内；
涡量为零的区域环量必为零——这正说明\textbf{无旋流动必无环量}（反之不必然，纯环流有环量却除涡心外无旋，
故 7.2 节的纯环流必须把涡心排除在积分回路之外）。\\
\textbf{教材式号与教材原话（书 p143--146）：}涡线方程为教材式 7.81；涡通量为教材式 7.82（$J=\displaystyle\iint_A 2\omega_n\,\dd A$），教材定义是“\textbf{涡通量是指通过任一曲面的旋转角速度矢量的通量的 2 倍}”；速度环量为教材式 7.83，教材规定“\textbf{平面闭曲线的正向规定为逆时针方向}，空间闭曲线的正向由需要确定”；斯托克斯定理为教材式 7.84，教材表述是“\textbf{沿有旋流场中一闭曲线由速度矢量产生的环量 $\Gamma$ 等于该闭曲线内所有涡通量之和}”。\\
\textbf{教材单独给出的复连通域情形（书 p146）：}“\textbf{通过复连通域的涡通量等于沿这个区域的外周线的速度环量与所有内周线的速度环量总和之差}”，即 $\Gamma_K-\sum\Gamma_k=\displaystyle\iint_A 2\omega_n\,\dd A$（教材图 7.29 用叶片叶型说明：把区域沿 AB 切开即化为单连通域）。\\
\textbf{教材 7.6 节开篇原话（书 p143）：}“在一流场中，如果一个区域内处处旋转角速度矢量 $\bomega$ 都为 0，这一区域内的流动是无旋的即有势的。理想流体的流动可以是有势的，也可以是有旋的。但\textbf{黏性流体的流动一般是有旋的}”；“\textbf{涡线是有旋流场中某一瞬时的一条曲线，曲线上各点处的旋转角速度矢量都与这一曲线相切}”；教材并指出“恒定流动中涡线形状不随时间变化”“与涡管垂直的断面称为涡管断面，微小断面的涡管称为微元涡管。涡管内充满着做旋转运动的流体称为涡束”。
\end{gongshi}

\begin{zhenti}{教材【例 7.4】（书 p146）}{证明题}
试证明均匀流的速度环量等于零。
\end{zhenti}
\noindent\textbf{教材证法：}① 取矩形封闭曲线（教材图 7.30a）：$\Gamma=\Gamma_{12}+\Gamma_{23}+\Gamma_{34}+\Gamma_{41}=bu+0-bu+0=0$；
② 再取圆周线（教材图 7.30b）：$\Gamma_K=\displaystyle\int_0^{2\pi}u\sin\theta\cdot r\,\dd\theta=ur\int_0^{2\pi}\sin\theta\,\dd\theta=0$；
③ 教材结论原话：“\textbf{同样，可以证明均匀流沿其他任何形状的封闭曲线的速度环量等于零}”。

\textbf{常考题型：}选择（"涡量为零"，见下）；填空（"速度环量 $\Gamma$ 的方向：逆时针为\rule{1.2cm}{0.4pt}"$\to$ 正）。
\textbf{重要度 \xstarfull{3}}\quad\yiju{E5 题库 \texttt{10.pdf} 选择第 13 题（"流体作无旋运动的特征"答"任意一点的涡量都为零"）；E2"掌握基本概念"}\quad\nandu{易}

\begin{zhenti}{题库 \texttt{10.pdf} 选择第 13 题}{单项选择题}
流体作无旋运动的特征是（\quad）。\\
(A) 所有流线都是直线\quad (B) 所有迹线都是直线\\
(C) 任意流体元的角变形为零\quad (D) 任意一点的涡量都为零
\end{zhenti}
\noindent\textbf{答案：(D)。}无旋的判据只看\textbf{微团自身的旋转}（涡量为零），
与流线、迹线是不是直线无关（同心圆流动可以是无旋的，直线流动也可以是有旋的）；
(C) 说的是"角变形"，那是有旋与否之外的另一个量（切应力相关），不能混为一谈。

\begin{yicuo}
\textbf{（本章第五号易错点）三个"为零"混淆：}\\
\textbf{涡量为零}（$\bomega=0$）＝无旋流动，这是无旋的定义；
\textbf{角变形速度为零}＝流体微团不变形（不是无旋）；
\textbf{流线是直线}与是否无旋\textbf{无关}。
2015 年选择第 4 题也考同一个点（答"微团无旋转的流动"），
"同心圆周运动"与"无旋"并不矛盾——微团可以不转着沿圆轨走（纯环流）。
\end{yicuo}
\end{kaodian}

\begin{kaodian}{3}{有旋流（旋涡流动）的压强分布}
\begin{gongshi}{旋涡中沿半径方向的压强分布}
\[ \frac{\dd p}{\dd r}=\rho\frac{u^2}{r} \]
\textbf{含义：}在有旋流动（如绕竖直轴的旋转容器内流体）中，沿半径方向取微元体列径向平衡：
外侧压强大于内侧，其差恰好提供流体的\textbf{离心惯性力}，
于是得到 $\dd p/\dd r=\rho u^2/r>0$——\textbf{外缘压强高于中心压强}。
积分（对刚性式旋转 $u=\omega r$）得 $p=p_0+\frac12\rho\omega^2r^2$，自由面为旋转抛物面。\\
\textbf{适用条件：}理想或黏性流体作定常轴对称旋转、忽略质量力沿径向的分量；
$u$ 为该半径处的切向速度。
\end{gongshi}
\textbf{常考题型：}选择/填空（"旋涡中心压强比外缘\rule{1.2cm}{0.4pt}"$\to$ 低）；
与第 2 章"相对平衡"的等角速度旋转容器互相印证，可作小计算。
\textbf{重要度 \xstarfull{3}}\quad\yiju{E2"掌握基本概念＋公式套用即可"；E6 教材 7.6 节；第 2 章相对平衡同型公式}\quad\nandu{易}
\end{kaodian}

\section{理想流体旋涡运动的基本定理}

\begin{kaodian}{3}{汤姆孙（开尔文）定理}
\begin{gongshi}{汤姆孙（开尔文）环量守恒定理（教材 7.7.1 节，书 p147，教材未给式号）}
\[ \Gamma=\oint_L \bu\cdot\dd\bl=\text{常数}\qquad\text{（沿同一条封闭流体线，不随时间变化）} \]
\textbf{含义：}在\textbf{理想、正压}流体中，当\textbf{质量力有势}时，
沿任意一条封闭\textbf{流体线}（由同一批流体质点组成的封闭曲线）的速度环量
\textbf{不随时间变化}——环量守恒。\\
\textbf{适用条件（三项缺一不可）：}① 理想流体（无黏性）；② 正压流体（密度只是压强的函数，
$\rho=\rho(p)$，均质不可压缩流体满足）；③ 质量力有势（如重力）。
三者任一不满足，环量都可能随时间变化。\\
\textbf{教材原话（书 p146--147）：}定理本体是“\textbf{汤姆逊定理可叙述为理想的不可压缩流体在有势质量力作用下沿一封闭流线体的速度环量不随时间变化}”，教材接着说明“作用于流体的重力是一有势力，因而在重力场中的理想正压流体将满足汤姆逊定理”——即\textbf{本体按“理想不可压缩流体＋有势质量力”叙述，再推广到正压流体}。\\
\textbf{教材两个必背定义（书 p146）：}“\textbf{流线体是指由相同流体质点组成的线状体。流线体随主流一起运动，流动中其位置与形状均可能发生变化，但构成它的流体质点始终不变}”；“\textbf{正压流体指密度仅随压强变化的流体，液体可以视为正压流体}”。\\
\textbf{（对旧稿的订正·术语与适用条件口径）}旧稿把封闭曲线写成“封闭流体线”，教材的专门术语是\textbf{流线体}；旧稿把三条条件并列成“理想＋正压＋有势”，而教材 7.7.1 节本体是“\textbf{理想不可压缩流体}＋有势质量力”，到 7.7.2 节才用“\textbf{正压}的理想流体在重力作用下”。\textbf{两种写法都能得分，但填空、简答按教材口径写更稳}。\\
\textbf{推论（教材原话，书 p147）：}“\textbf{理想流体如果开始作无旋流动，流动将永远是有势的}”；
反之，若某时刻流场无旋（处处 $\Gamma=0$），则涡线不会自行产生，有旋也不会凭空消失。
教材的推理链是：静止理想流体沿闭曲线的环量为零 $\to$ 由汤姆逊定理环量恒为零 $\to$ 由斯托克斯定理涡通量恒为零 $\to$ 流动始终无旋。
\end{gongshi}
\textbf{常考题型：}选择/判断（"理想正压流体、质量力有势时，沿封闭曲线的速度环量\rule{1.6cm}{0.4pt}"$\to$ 守恒／不随时间变化）；
简答（写出定理及三个适用条件）。
\textbf{重要度 \xstarfull{3}}\quad\yiju{E2"第六到第八章只有少量填空/判断/选择"；E1 slide33"概念为主"；E6 教材 7.7 节}\quad\nandu{中}
\end{kaodian}

\begin{kaodian}{3}{亥姆霍兹三定理与拉格朗日涡定理}
\begin{gongshi}{亥姆霍兹三定理（涡管守恒三条，教材 7.7.2 节，书 p147--148）}
\begin{enumerate}
  \item \textbf{亥姆霍兹第一定理：}“在同一时刻，通过涡管任意涡管断面的涡通量不变”（涡管强度沿管不变，$\Omega_1\dd A_1=\Omega_2\dd A_2$）；
  \item \textbf{亥姆霍兹第二定理：}“正压的理想流体在重力作用下，组成涡管的流体质点将始终组成涡管”（涡管随流运动时位置与形状可变，但管内流体质点不变）；
  \item \textbf{亥姆霍兹第三定理：}“正压的理想流体在重力作用下，涡管强度不随时间而变化”。
\end{enumerate}
\textbf{教材原话（书 p147--148）：}上面三条就是教材原文，直接背这三句。\\
\textbf{适用条件：}第二、第三定理的教材表述都含“\textbf{正压的理想流体在重力作用下}”，与汤姆逊定理一致；
教材另加一句提醒（书 p148）：“在使用上述定理时应注意它们的应用条件。工程中实际不存在理想流体，但在短时间内，可以认为流体运动满足这些条件，从而简化研究过程。”
\end{gongshi}

\begin{tishi}
\textbf{（对旧稿的订正·亥姆霍兹三定理错位）}旧稿把“亥姆霍兹三定理”写成① 涡管强度不变、② 涡管不能中断、③ 涡线随流体运动，\textbf{与教材对不上}，差别有二：\\
① 教材的\textbf{第二定理是“组成涡管的流体质点将始终组成涡管”}（＝旧稿的 ③），而教材的\textbf{第三定理是“涡管强度不随时间而变化”——旧稿整条缺失}，必须补上；\\
② 旧稿的“涡管不能中断”\textbf{在教材中不是独立定理}，而是第一定理的\textbf{推论}。教材原话（书 p147）：“涡管断面面积不能减小到 $0$，否则这里的旋转角速度矢量的大小将趋近于无穷大，而这是不可能的。因此，流场中的\textbf{涡管首尾断面只能终止于液面或固壁处，或者涡管成为环状}。”选择题问“涡管能否在流体内部中断”答“不能”，依据写这句推论即可。\\
（旧稿的 ② 不是错误结论，只是\textbf{不该占“定理”的位置}；现按教材把三条定理写准，推论另列。）
\end{tishi}

\begin{gongshi}{拉格朗日涡定理}
\[ \text{理想正压流体、质量力有势，若某时刻流动无旋，则永远无旋} \]
\textbf{含义：}常密度理想流体在重力（有势质量力）作用下，若初始时刻流场无旋
（即处处 $\bomega=0$，涡通量全为零），则以后时刻流场始终保持无旋。
它是汤姆孙定理的推论：无旋时任意封闭流体线的环量恒为零，因而涡通量恒为零。\\
\textbf{适用条件：}理想、正压（均质不可压缩）、质量力有势——与汤姆孙定理完全一致。
\textbf{工程意义：}这正是可以把远离固体壁面、绕过物体的外部流动按"势流"处理的根据：
来流无旋、流体近似理想，则绕流区无旋。
\end{gongshi}

\textbf{常考题型：}选择/判断（"涡管在流体中能否中断"$\to$ 不能；"理想流体中无旋流动恒为无旋"）；
填空（"亥姆霍兹三定理的内容：\rule{2cm}{0.4pt}"）。
\textbf{重要度 \xstarfull{3}}\quad\yiju{E2"第六到第八章只有少量填空/判断/选择，掌握基本概念即可"；E6 教材 7.7 节}\quad\nandu{中}

\begin{tishi}
四条定理的\textbf{适用条件完全一样}（理想、正压、质量力有势），所以做题时它们经常一起出现；
只要题目里\emph{没有任何一条}提到黏性、提到斜压、提到质量力无势，就可以放心用。
反过来，题目若问"哪一条不满足时定理失效"，答"黏性／斜压／质量力无势"三者之一即可。
\end{tishi}
\end{kaodian}

\begin{kaodian}{3}{毕奥-萨伐尔定理（涡对诱导速度）}
\begin{gongshi}{毕奥-萨伐尔定理（教材 7.8 节，书 p148）}
\[ \dd\bv=\frac{\Gamma}{4\pi}\frac{\dd\bl\times\br}{r^3} \]
\textbf{含义：}一条强度为 $\Gamma$ 的\textbf{涡线段} $\dd\bl$ 在其周围空间点（相对位置矢量 $\br$、
距离 $r$）诱导出的速度 $\dd\bv$，\textbf{教材原话（书 p148）：}“计算诱导速度可利用涡线与通电导线之间的相似关系，认为\textbf{涡线相当于通电导线，涡线周围产生的速度场 $u$ 相当于通电导线在其周围感应产生的磁场 $B$}”；
“\textbf{诱导速度的方向由右手法则确定：大拇指表示旋涡转动方向，则诱导速度的方向由其余四指给出}”（教材把电流磁场公式 $\dd B=\dfrac{I\sin\alpha\,\dd l}{4\pi r^2}$ 与涡束对照，导入诱导速度）。
旧稿写的“四指顺环量、拇指指涡线走向”方向与此一致，但考场若考“右手定则怎么用”，按教材这句写最稳。
整个涡线（涡管）的诱导速度是各微段的叠加。\\
\textbf{适用条件：}不可压缩流体；涡线在流场内部（公式给出的是该涡线在其他点产生的诱导速度，
在涡线上本身 $r\to0$ 发散，需另行处理）。它把"涡"与"速度场"联系起来：
\textbf{已知涡的分布 $\to$ 求速度场}，是涡动力学与升力面理论的基础。\\
\textbf{常用结论：}无限长直线涡（强度 $\Gamma$）在距离 $r$ 处诱导的速度
$u=\dfrac{\Gamma}{2\pi r}$——正是 7.2 节纯环流的速度公式（二者互为印证）。\\
\textbf{教材式号（书 p148--149）：}任意形状涡束对任意点 $P$ 的诱导速度为教材式 7.85（$u=\displaystyle\int_l\frac{\Gamma\sin\theta}{4\pi r^2}\,\dd l$）；
\textbf{有限长直线涡束}为教材式 7.86（$u=\dfrac{\Gamma}{4\pi r_0}(\cos\theta_2-\cos\theta_1)$，见下框的 OCR 说明）；
\textbf{半无限长涡束} $u=\dfrac{\Gamma}{4\pi r_0}$、\textbf{无限长涡束} $u=\dfrac{\Gamma}{2\pi r_0}$（教材书 p149）——考试真正要记的是这两个特例。
\end{gongshi}

\begin{tishi}
\textbf{（OCR 判读存疑，两说并列）}教材式 7.86 的余弦次序，在扫描 OCR 中只能读到“(cos$\theta_2$--cos$\theta_1$)”这一残形，
而教材图 7.34 中 $\theta_1$、$\theta_2$ 标注在哪一端无法从 OCR 判读；若按 $\displaystyle\int_{\theta_1}^{\theta_2}\sin\theta\,\dd\theta$ 推，结果是 $\cos\theta_1-\cos\theta_2$。
\textbf{两说并列，不要擅自删原说法}：题目直接给教材式 7.86 就按 $\cos\theta_2-\cos\theta_1$，自己推积分就按 $\cos\theta_1-\cos\theta_2$（两者只差一个符号，取决于 $\theta$ 从哪一端量起）。
好在\textbf{半无限长与无限长两个特例不受次序影响}，考试优先记这两个。
\end{tishi}
\textbf{常考题型：}选择/填空（"无限长直线涡诱导速度与距离成\rule{1.2cm}{0.4pt}比"$\to$ 反比；
"诱导速度的方向由\rule{1.6cm}{0.4pt}定则确定"$\to$ 右手）。
\textbf{重要度 \xstarfull{3}}\quad\yiju{E2"掌握基本概念＋公式套用"；E6 教材 7.7、7.8 节}\quad\nandu{中}
\end{kaodian}

\section{旋涡诱导速度}

\begin{kaodian}{1}{直线涡、涡对与涡列的诱导速度}
\begin{gongshi}{单根无限长直线涡（教材式 7.85、7.86 的特例）}
\[ u=\frac{\Gamma}{2\pi r} \]
\textbf{含义：}强度为 $\Gamma$ 的无限长直线涡，在距涡线 $r$ 处诱导的速度大小与 $r$ 成反比；
速度方向垂直于"涡线—场点"连线的平面，由右手定则定出。\\
\textbf{适用条件：}无限长直涡线、不可压缩流体、场点不在涡线上。
\end{gongshi}

\begin{gongshi}{涡对与涡列（简介；教材 7.8.2 节只讲平面涡层）}
\[ \text{等强度反向平行涡对（涡对）：两涡连线的中垂线上速度最大} \]
\[ \text{同向涡列（卡门涡街模型）：涡列以一定速度向下游平移} \]
\textbf{含义：}两根平行涡线组成\textbf{涡对}：\\等强度同向时为"涡旋对"，
两涡绕公共中心互相诱导旋转，如同刚体式旋转；等强度反向时两涡一起以相同速度平移，
不会彼此靠近或远离（工程中用它模拟翼尖涡）。多根涡线平行排列构成\textbf{涡列}：
各涡互相诱导，使整列涡以某个确定的\textbf{平移速度}沿垂直于涡线的方向移动；
这一结论是解释圆柱绕流"卡门涡街"尾涡形式的基础。\\
\textbf{适用条件：}不可压缩流体、无限长直涡线、各涡强度恒定（无黏）；
涡列平移速度由各涡的诱导速度叠加得到。\\
\textbf{教材原话（书 p150）：}“\textbf{在无限流场中布置一涡列，这一涡列由多个无限长涡束无间隔地直线排列而成，称为涡层}”；“由于 $\gamma$ 为负值，在涡层上方其速度为正，下方为负，即\textbf{当穿过涡层时，切向速度将产生间断（跃变）$\gamma$，而法向速度则是连续的}”（教材 7.8.2 节）。\\
\textbf{教材覆盖说明：}教材 7.8.2 节只讨论\textbf{平面涡层}的速度跃变；本章已 OCR 的 31 页正文中\textbf{没有出现“涡对”}，也没有“两涡互相诱导绕转、涡列以确定速度向下游平移”的内容。
本框的涡对、涡列结论来自外校辅导书整理稿，真题中未见成题，故重要度保持 \xstarfull{1}；考试认\textbf{“涡层切向速度间断、法向速度连续”}这句教材原话即可。
\end{gongshi}
\textbf{常考题型：}填空/选择（涡对、涡列的诱导速度方向判断）——真题中未见，属了解性内容。
\textbf{重要度 \xstarfull{1}}\quad\yiju{E4 历年真题中未见成题；E1 slide33"第七章＝概念为主"（仅章节级依据）}\quad\nandu{难}
\end{kaodian}

\section{本章小结}

\begin{center}
\begin{tabularx}{\textwidth}{@{}XccX@{}}
\toprule
考点 & 重要度 & 难度 & 一句话抓住 \\
\midrule
速度势函数 $\varphi$ 的定义与存在条件 & \xstarfull{5} & 中 & 无旋才有 $\varphi$；$\bu=\grad\varphi$，不可压缩时 $\grad^2\varphi=0$ \\
流函数 $\psi$ 的定义与流量关系 & \xstarfull{5} & 中 & 平面就有 $\psi$；$\psi=c$ 是流线，$\Delta\psi$＝单宽流量 \\
$\varphi$ 与 $\psi$ 的关系（正交、CR、互求） & \xstarfull{5} & 中 & 处处正交；$\varphi_x=\psi_y$、$\varphi_y=-\psi_x$ \\
均匀直线流 & \xstarfull{4} & 易 & $\varphi=Ux$、$\psi=Uy$ \\
点源与点汇 & \xstarfull{4} & 易 & $u_r=Q/(2\pi r)$；$\psi=Q\theta/2\pi$ \\
偶极子 & \xstarfull{4} & 中 & 源汇无限接近的极限；$M=2aq$，$M$ 由汇指向源 \\
教材 7.1.4--7.1.5 节 复势与复速度 & \xstarfull{2} & 中 & $W=\varphi+\mathrm{i}\psi$；$V=\dd W/\dd z=u_x-\mathrm{i}u_y$ \\
教材 7.3.2 节 螺旋流（点汇＋点涡） & \xstarfull{2} & 中 & 等势线与流线都是对数螺旋线 $r=C_1\mathrm{e}^{q\theta/\Gamma}$ \\
纯环流（势涡） & \xstarfull{4} & 中 & 有势无旋；$u_\theta=\Gamma/(2\pi r)$ \\
势流叠加、驻点与兰肯半体 & \xstarfull{4} & 中 & $\varphi=\sum\varphi_i$；驻点令 $u_x=u_y=0$ \\
圆柱体绕流（无环量） & \xstarfull{3} & 中 & 均匀流＋偶极子；$|u_\theta|_{\max}=2U$ \\
达朗贝尔佯谬与库塔-儒可夫斯基升力 & \xstarfull{3} & 中 & 无环量阻力为零；升力 $L=\rho U\Gamma$ \\
空间势流与附加质量 & \xstarfull{1} & 中 & $u_r=Q/(4\pi r^2)$；附加质量＝带动流体的惯性效应 \\
涡线、涡管、涡通量与速度环量 & \xstarfull{3} & 易 & 环量＝涡通量（斯托克斯定理） \\
有旋流的压强分布 & \xstarfull{3} & 易 & $\dd p/\dd r=\rho u^2/r$，外缘高、中心低 \\
汤姆孙定理与亥姆霍兹三定理 & \xstarfull{3} & 中 & 理想＋正压＋质量力有势 $\to$ 环量守恒 \\
拉格朗日涡定理 & \xstarfull{3} & 中 & 无旋恒无旋（势流处理的根据） \\
毕奥-萨伐尔定理 & \xstarfull{3} & 中 & $\dd\bv=\frac{\Gamma}{4\pi}\frac{\dd\bl\times\br}{r^3}$；长直涡 $u=\Gamma/2\pi r$ \\
旋涡诱导速度（直线涡、涡对、涡列） & \xstarfull{1} & 难 & 直线涡 $u=\Gamma/2\pi r$；涡对反向平移 \\
\bottomrule
\end{tabularx}
\end{center}
